\documentclass[12pt]{article}

\usepackage[a4paper,margin=1in]{geometry}
\usepackage{setspace}
\onehalfspacing

\usepackage{graphicx}
\usepackage{amsmath,amssymb,amsfonts,amsthm}
\usepackage{bm}
\usepackage{mathtools}
\usepackage{booktabs}
\usepackage{enumitem}
\usepackage{hyperref}

\setlength{\parindent}{0pt}
\setlength{\parskip}{0.6em}

\setlength{\abovedisplayskip}{10pt}
\setlength{\belowdisplayskip}{10pt}
\setlength{\abovedisplayshortskip}{6pt}
\setlength{\belowdisplayshortskip}{6pt}

\newtheorem{proposition}{Proposition}
\newtheorem{remark}{Remark}

\newcommand{\E}{\mathbb{E}}
\newcommand{\KL}{\operatorname{KL}}
\newcommand{\cov}{\operatorname{cov}}
\newcommand{\Var}{\operatorname{Var}}
\newcommand{\tr}{\operatorname{tr}}
\newcommand{\diag}{\operatorname{diag}}
\newcommand{\vecop}{\operatorname{vec}}

\title{Annotated Joint Online Gaussian Process with Linear Mean Model\\
for Kronecker-Separable Spatio-Temporal HiPPO-SVGP}
\author{}
\date{March 2026\\Annotated English explanation version}

\begin{document}
\maketitle

\section{Joint Observation Model}

We consider a spatio-temporal latent response observed at spatial locations
$s_i\in\mathcal{S}\subset\mathbb{R}^{d_s}$ and time points $t_n\in\mathbb{R}$.
The observation at location $(t_n,s_i)$ is modeled as the sum of a linear mean
component, a nonparametric Gaussian process component, and independent Gaussian noise:
\begin{equation}
y(t_n,s_i)
=
\phi(t_n,s_i)^\top \beta
+
f(t_n,s_i)
+
\epsilon_{n,i},
\qquad
\epsilon_{n,i}\sim\mathcal{N}(0,\sigma^2).
\label{eq:joint-observation-scalar}
\end{equation}
Here $\phi(t_n,s_i)\in\mathbb{R}^{d}$ is a covariate design vector,
$\beta\in\mathbb{R}^{d}$ is the linear mean parameter, and
$f$ is a spatio-temporal Gaussian process.

Let the data in the $n$-th online temporal block be
\begin{equation}
D_n=\{(t,s,y,\phi): (t,s)\in\mathcal{B}_n\},
\end{equation}
and let
\begin{equation}
y_n\in\mathbb{R}^{B_n},
\qquad
\Phi_n\in\mathbb{R}^{B_n\times d}
\end{equation}
denote the stacked response vector and design matrix for this block. The block-level
observation model is
\begin{equation}
y_n = \Phi_n \beta + f_n + \epsilon_n,
\qquad
\epsilon_n\sim\mathcal{N}(0,\sigma^2 I).
\label{eq:block-observation}
\end{equation}

The Gaussian process prior is assumed to be separable:
\begin{equation}
f\sim \mathcal{GP}(0,k),
\qquad
k\big((s,t),(s',t')\big)
=
k_s(s,s')k_t(t,t').
\label{eq:separable-kernel}
\end{equation}
If latent values are vectorized over a full time--space grid, then
\begin{equation}
K_{ff}=K_t\otimes K_s.
\label{eq:full-kron-prior}
\end{equation}

The key modeling decision in this document is to train the linear mean and the GP
component jointly. Instead of first estimating $\beta$ and then fitting a GP to residuals,
we directly construct an online posterior over the joint unknowns $(\beta,u)$, where
$u$ denotes the GP inducing variables.

\section{Kronecker HiPPO-SVGP Inducing Representation}

\subsection{Temporal HiPPO-RFF Interdomain Features}

For the temporal axis, we use HiPPO-LegS interdomain inducing variables together with
a random Fourier feature approximation of a stationary temporal kernel.

For a time horizon $T>0$, define the scaled Legendre basis
\begin{equation}
g^{(T)}_\ell(x)
=
\sqrt{\frac{2\ell+1}{T}}\,
P_\ell\!\left(\frac{2x}{T}-1\right),
\qquad
x\in[0,T],
\quad
\ell=0,\dots,M_t-1,
\label{eq:scaled-legs-basis}
\end{equation}
where $P_\ell$ is the Legendre polynomial of degree $\ell$.
The temporal interdomain inducing variables are
\begin{equation}
u_\ell^{(t)}(T)
=
\int_0^T g_\ell^{(T)}(x) f_t(x)\,dx.
\label{eq:temporal-hippo-inducing}
\end{equation}

Assume the temporal kernel $k_t$ is stationary. By Bochner's theorem, we approximate it
using random Fourier features:
\begin{equation}
k_t(\tau,\tau')
\approx
\sum_{r=1}^{R} w_r
\left[
\cos(\omega_r\tau)\cos(\omega_r\tau')
+
\sin(\omega_r\tau)\sin(\omega_r\tau')
\right].
\label{eq:rff-kernel}
\end{equation}
Define the oscillatory integrals
\begin{align}
I^{\cos}_\ell(T;\omega)
&=
\int_0^T \cos(\omega x)g_\ell^{(T)}(x)\,dx,\\
I^{\sin}_\ell(T;\omega)
&=
\int_0^T \sin(\omega x)g_\ell^{(T)}(x)\,dx.
\end{align}
Using the spherical-Bessel identity
\begin{equation}
\int_{-1}^{1} e^{i\kappa u}P_\ell(u)\,du
=
2 i^\ell j_\ell(\kappa),
\qquad
\kappa=\frac{\omega T}{2},
\end{equation}
these integrals admit closed-form expressions involving spherical Bessel functions
$j_\ell(\cdot)$. This gives analytic temporal inducing covariances and point-to-inducing
cross-covariances.

Let $B_t(T)\in\mathbb{R}^{M_t\times 2R}$ be the temporal HiPPO-RFF feature matrix,
with entries
\begin{equation}
[B_t(T)]_{\ell,r}
=
\sqrt{w_r}\,I^{\cos}_\ell(T;\omega_r),
\qquad
[B_t(T)]_{\ell,R+r}
=
\sqrt{w_r}\,I^{\sin}_\ell(T;\omega_r),
\end{equation}
and define the temporal RFF vector
\begin{equation}
\varphi_t(\tau)
=
\left(
\sqrt{w_1}\cos(\omega_1\tau),\dots,\sqrt{w_R}\cos(\omega_R\tau),
\sqrt{w_1}\sin(\omega_1\tau),\dots,\sqrt{w_R}\sin(\omega_R\tau)
\right)^\top.
\end{equation}
Then
\begin{equation}
K_{uu}^{(t)}(T)
\approx
B_t(T)B_t(T)^\top,
\label{eq:temporal-inducing-cov}
\end{equation}
and
\begin{equation}
k_{fu}^{(t)}(\tau;T)
\approx
\varphi_t(\tau)^\top B_t(T)^\top.
\label{eq:temporal-cross-cov}
\end{equation}

For local temporal mini-batches, the same construction is applied using a local horizon
$T_n$ for block $D_n$, yielding $K_{uu,n}^{(t)}$ and $K_{fu,n}^{(t)}$.

\subsection{Spatial Inducing Variables}

On the spatial axis, we use inducing locations
\begin{equation}
Z_s=\{z_1^{(s)},\dots,z_{M_s}^{(s)}\}.
\end{equation}
The spatial inducing covariance and cross-covariance are
\begin{align}
K_{ZZ}^{(s)}
&\in\mathbb{R}^{M_s\times M_s},
&
[K_{ZZ}^{(s)}]_{jj'}
&=
k_s(z_j^{(s)},z_{j'}^{(s)}),
\\
K_{XZ}^{(s)}
&\in\mathbb{R}^{N_s\times M_s},
&
[K_{XZ}^{(s)}]_{ij}
&=
k_s(s_i,z_j^{(s)}).
\end{align}

The spatial inducing set is fixed across online time steps. This is important: online
adaptation occurs through the temporal HiPPO component, while the spatial inducing
geometry remains shared across all blocks.

\subsection{Mixed Spatio-Temporal Inducing Variables}

For each spatial inducing location $z_j^{(s)}$, define temporal interdomain variables
\begin{equation}
u_{\ell,j}(T)
=
\int_0^T g_\ell^{(T)}(x)f(x,z_j^{(s)})\,dx.
\end{equation}
Stacking the temporal and spatial inducing variables gives
\begin{equation}
u(T)
=
\vecop\!\left([u_{\ell,j}(T)]_{\ell,j}\right)
\in\mathbb{R}^{M_tM_s}.
\end{equation}

By separability of the spatio-temporal kernel,
\begin{equation}
K_{uu}^{(st)}(T)
=
K_{uu}^{(t)}(T)
\otimes
K_{ZZ}^{(s)}.
\label{eq:mixed-inducing-cov}
\end{equation}
For a block $D_n$, the point-to-inducing covariance has the corresponding Kronecker
form
\begin{equation}
K_{fu,n}^{(st)}
=
K_{fu,n}^{(t)}
\otimes
K_{XZ}^{(s)}.
\label{eq:mixed-cross-cov}
\end{equation}

The sparse GP projection matrix for block $n$ is therefore
\begin{equation}
A_n
=
K_{fu,n}^{(st)}
\left(K_{uu,n}^{(st)}\right)^{-1}.
\label{eq:projection-general}
\end{equation}
Using Kronecker identities,
\begin{equation}
A_n
=
\left[
K_{fu,n}^{(t)}
\left(K_{uu,n}^{(t)}\right)^{-1}
\right]
\otimes
\left[
K_{XZ}^{(s)}
\left(K_{ZZ}^{(s)}\right)^{-1}
\right].
\label{eq:projection-kron}
\end{equation}
Define
\begin{equation}
T_n
:=
K_{fu,n}^{(t)}
\left(K_{uu,n}^{(t)}\right)^{-1},
\qquad
C
:=
K_{XZ}^{(s)}
\left(K_{ZZ}^{(s)}\right)^{-1}.
\label{eq:Tn-C-def}
\end{equation}
Then
\begin{equation}
A_n=T_n\otimes C.
\label{eq:A-kron}
\end{equation}

\section{Why Residual-Based Online Learning is Inconsistent}

A two-stage procedure first estimates the linear mean and then fits a GP to residuals:
\begin{equation}
r_n(\hat{\beta}_n)
=
y_n-\Phi_n\hat{\beta}_n.
\end{equation}
The difficulty is that residuals are not fixed targets. If the estimate of the linear
coefficient changes from $\hat{\beta}_{n-1}$ to $\hat{\beta}_{n}$, then for any historical
block $j<n$,
\begin{align}
r_j(\hat{\beta}_{n})
&=
y_j-\Phi_j\hat{\beta}_{n}
\\
&=
y_j-\Phi_j\hat{\beta}_{n-1}
-
\Phi_j(\hat{\beta}_{n}-\hat{\beta}_{n-1})
\\
&=
r_j(\hat{\beta}_{n-1})
-
\Phi_j(\hat{\beta}_{n}-\hat{\beta}_{n-1}).
\label{eq:residual-mismatch}
\end{align}
Thus, the historical residual memory becomes stale whenever $\beta$ is updated.

\begin{proposition}[Joint training removes stale residual targets]
In the joint model
\begin{equation}
y_j\mid \beta,u_j
\sim
\mathcal{N}(\Phi_j\beta+A_j u_j,\sigma^2I),
\end{equation}
historical information is stored as a likelihood contribution to the same joint posterior
over $(\beta,u)$, rather than as residual observations computed from a changing point
estimate of $\beta$. Therefore, historical information does not become inconsistent when
the posterior mean of $\beta$ changes.
\end{proposition}

\begin{proof}
The two-stage target for block $j$ is
$r_j(\hat{\beta})=y_j-\Phi_j\hat{\beta}$, which explicitly depends on the current point
estimate $\hat{\beta}$. Equation~\eqref{eq:residual-mismatch} shows that the target
changes when $\hat{\beta}$ changes.

By contrast, joint training uses the likelihood
\begin{equation}
p(y_j\mid \beta,u_j)
=
\mathcal{N}(y_j\mid \Phi_j\beta+A_j u_j,\sigma^2I).
\end{equation}
This likelihood is a fixed probabilistic statement about $y_j$ as a function of the
unknown variables $(\beta,u_j)$. Updating the posterior mean of $\beta$ changes our
belief about $\beta$, but it does not redefine the historical observation $y_j$ or the
likelihood contribution of block $j$. Hence no stale residual target is created.
\end{proof}

The essential distinction is
\begin{equation}
\text{two-stage: learn } y-\Phi\hat{\beta},
\qquad
\text{joint training: learn } p(y\mid \beta,u).
\end{equation}
This is the reason the rest of this document focuses on a joint online posterior over
the linear mean and GP inducing variables.


\section{Online Structured Joint Training ELBO}
\label{sec:online_structured_joint_elbo}

For online block $n$, the sparse GP approximation gives
\begin{equation}
    f_n\mid u_n \approx A_nu_n .
\end{equation}
The likelihood is therefore
\begin{equation}
    p(y_n\mid \beta,u_n)
    =
    \mathcal{N}
    \left(
        y_n\mid \Phi_n\beta+A_nu_n,\sigma^2I
    \right).
    \label{eq:joint-likelihood}
\end{equation}

The main method in this document uses a \emph{structured joint Gaussian posterior}
over the linear coefficient and the GP inducing variables:
\begin{equation}
    q_n(\beta,u_n)
    =
    \mathcal{N}
    \left(
    \begin{bmatrix}\beta\\u_n\end{bmatrix}
    \middle|
    \begin{bmatrix}m_{\beta,n}\\m_{u,n}\end{bmatrix},
    S_n
    \right),
    \qquad
    S_n^{-1}=\Lambda_n .
    \label{eq:structured-joint-posterior}
\end{equation}
The precision is stored in block form
\begin{equation}
    \Lambda_n
    =
    \begin{bmatrix}
        \Lambda_{\beta\beta,n} & \Lambda_{\beta u,n}\\
        \Lambda_{u\beta,n} & \Lambda_{uu,n}
    \end{bmatrix},
    \qquad
    h_n=
    \begin{bmatrix}
        h_{\beta,n}\\h_{u,n}
    \end{bmatrix}.
    \label{eq:structured-joint-precision}
\end{equation}
The crucial modeling choice is that $\Lambda_{uu,n}$ keeps the
Kronecker/Sylvester structure, while $\Lambda_{\beta u,n}$ is retained rather
than discarded. Since the dimension $d$ of $\beta$ is usually small, the
cross block $\Lambda_{\beta u,n}\in\mathbb{R}^{d\times M_tM_s}$ can be stored
explicitly or as $d$ reshaped temporal--spatial matrices. This preserves the
posterior dependence between the linear trend and the GP component without
requiring a dense covariance over all inducing variables.

Let the base prior for the static linear coefficient be
\begin{equation}
    p(\beta)=\mathcal{N}(m_{\beta,0},K_\beta),
    \qquad
    P_\beta:=K_\beta^{-1},
\end{equation}
and let the current GP inducing prior be
\begin{equation}
    p(u_n)=\mathcal{N}(0,K_{nn}).
\end{equation}
The online variational objective can be written as
\begin{equation}
    \mathcal{L}_n
    =
    \E_{q_n(\beta,u_n)}
    \left[
        \log p(y_n\mid \beta,u_n)
    \right]
    -
    \KL
    \left[
        q_n(\beta,u_n)
        \Vert
        p_n(\beta,u_n)
    \right],
    \label{eq:structured-joint-elbo}
\end{equation}
where $p_n(\beta,u_n)$ denotes the prior plus the transferred old-likelihood
information defined in Section~\ref{sec:changing_basis_transfer}. In the
Gaussian case this objective is equivalent to updating the joint natural
parameters.

For a full joint Gaussian posterior, the expected log likelihood is
\begin{align}
    \E\log p(y_n\mid \beta,u_n)
    &=-\frac{B_n}{2}\log(2\pi\sigma^2) \\
    &\quad
    -\frac{1}{2\sigma^2}
    \Big[
        \|y_n-\Phi_nm_{\beta,n}-A_nm_{u,n}\|^2
        +\tr(\Phi_nS_{\beta\beta,n}\Phi_n^\top)  \\
    &\qquad\qquad
        +\tr(A_nS_{uu,n}A_n^\top)
        +2\tr(\Phi_nS_{\beta u,n}A_n^\top)
    \Big].
    \label{eq:structured-expected-log-likelihood}
\end{align}
The cross term $2\tr(\Phi_nS_{\beta u,n}A_n^\top)$ is exactly the term lost by
a mean-field approximation $q(\beta)q(u_n)$. Retaining it is the reason for
using the structured joint route in the main method.

A fully factorized posterior remains useful as an ablation, but it is not the
main formulation. In the remainder of the document, ``main method'' refers to
the structured joint precision update in which $\Lambda_{\beta u,n}$ is kept
and posterior moments are recovered by a Schur-complement computation.

\section{Fixed-Basis Exact Joint Gaussian Update}

If the inducing representation is fixed across online blocks, i.e. the same inducing
variable $u$ is used throughout, then the joint model is conjugate under a Gaussian
likelihood. In this special case, the exact joint posterior can be maintained in
information form.

Define the joint state
\begin{equation}
z=
\begin{bmatrix}
\beta\\
u
\end{bmatrix},
\qquad
H_n=
\begin{bmatrix}
\Phi_n & A_n
\end{bmatrix}.
\label{eq:joint-state-fixed-basis}
\end{equation}
Then
\begin{equation}
y_n\mid z
\sim
\mathcal{N}(H_nz,\sigma^2I).
\end{equation}
Let
\begin{equation}
q_{n-1}(z)=\mathcal{N}(m_{n-1},S_{n-1}),
\end{equation}
and define the natural parameters
\begin{equation}
\Lambda_{n-1}=S_{n-1}^{-1},
\qquad
h_{n-1}=S_{n-1}^{-1}m_{n-1}.
\end{equation}
The exact Gaussian-conjugate update is
\begin{equation}
\Lambda_n
=
\Lambda_{n-1}
+
\frac{1}{\sigma^2}H_n^\top H_n,
\label{eq:joint-precision-update}
\end{equation}
\begin{equation}
h_n
=
h_{n-1}
+
\frac{1}{\sigma^2}H_n^\top y_n.
\label{eq:joint-info-update}
\end{equation}
The posterior moments are recovered by
\begin{equation}
S_n=\Lambda_n^{-1},
\qquad
m_n=S_nh_n.
\label{eq:joint-moment-recovery}
\end{equation}

Expanding the precision contribution gives
\begin{equation}
H_n^\top H_n
=
\begin{bmatrix}
\Phi_n^\top\Phi_n
&
\Phi_n^\top A_n
\\
A_n^\top\Phi_n
&
A_n^\top A_n
\end{bmatrix}.
\label{eq:joint-precision-blocks}
\end{equation}
Therefore,
\begin{equation}
\Lambda_n
=
\Lambda_{n-1}
+
\frac{1}{\sigma^2}
\begin{bmatrix}
\Phi_n^\top\Phi_n
&
\Phi_n^\top A_n
\\
A_n^\top\Phi_n
&
A_n^\top A_n
\end{bmatrix},
\end{equation}
and
\begin{equation}
h_n
=
h_{n-1}
+
\frac{1}{\sigma^2}
\begin{bmatrix}
\Phi_n^\top y_n\\
A_n^\top y_n
\end{bmatrix}.
\end{equation}

This update is stronger than a mean-field approximation because it keeps the posterior
cross-covariance
\begin{equation}
\cov(\beta,u)\neq 0.
\end{equation}
The fixed-basis exact joint Gaussian update should therefore be used as a baseline and
sanity check for more scalable structured or mean-field approximations.

\begin{remark}
Equations~\eqref{eq:joint-precision-update}--\eqref{eq:joint-moment-recovery} are valid
only when all historical likelihood contributions are expressed in the same inducing
coordinates. If the HiPPO temporal inducing basis changes over time, the previous
posterior must first be transferred to the new inducing coordinates before the next
likelihood update is applied.
\end{remark}


\section{Changing-Basis Structured Joint Old-Likelihood Transfer}
\label{sec:changing_basis_transfer}

In the HiPPO-SVGP setting, the temporal inducing basis depends on the current
horizon or local block. Thus the old inducing variable $u_o$ and the new
inducing variable $u_n$ are generally different random variables. For example,
\begin{equation}
    u_{\ell}^{old}
    =
    \int_0^{T_{n-1}} g_\ell^{(T_{n-1})}(x)f(x)\,dx,
    \qquad
    u_{\ell}^{new}
    =
    \int_0^{T_n} g_\ell^{(T_n)}(x)f(x)\,dx .
\end{equation}
Therefore the previous posterior cannot be directly reused as a posterior over
$u_n$. The old likelihood information must be transferred from the old
coordinates $(\beta,u_o)$ to the new coordinates $(\beta,u_n)$.

\subsection{Streaming VFE Transfer for the Joint State}
\label{subsec:streaming_vfe_transfer}

The streaming sparse GP principle represents historical observations by the
old likelihood ratio. For the joint state, the old approximate likelihood
factor is
\begin{equation}
    \frac{q_{n-1}(\beta,u_o)}{p(\beta)p(u_o)} .
\end{equation}
The changing-basis VFE update becomes
\begin{equation}
    q_n(\beta,u_n)
    \propto
    p(\beta)p(u_n)
    \exp
    \left\{
    \E_{p(u_o\mid u_n)}
    \left[
        \log
        \frac{q_{n-1}(\beta,u_o)}{p(\beta)p(u_o)}
    \right]
    +
    \log p(y_n\mid \beta,u_n)
    \right\}.
    \label{eq:joint_streaming_vfe}
\end{equation}
Here $\beta$ is unchanged across the transfer, while $u_o$ is integrated through
the GP conditional $p(u_o\mid u_n)$. This is the structured joint analogue of
the old-likelihood-ratio transfer used in streaming sparse GP inference. It is
stronger than transferring only a marginal posterior over $u$, because the
historical correlation between the linear trend and the GP residual is retained.

\subsection{Closed-Form Gaussian Transfer of the Joint Old Likelihood}
\label{subsec:closed_form_old_likelihood_transfer}

Assume that the old joint posterior and the corresponding old prior are
Gaussian. Define the old joint likelihood natural parameters by
\begin{equation}
    \log
    \frac{q_{n-1}(\beta,u_o)}{p(\beta)p(u_o)}
    =
    -\frac{1}{2}
    \begin{bmatrix}\beta\\u_o\end{bmatrix}^{\!\top}
    R_o^z
    \begin{bmatrix}\beta\\u_o\end{bmatrix}
    +
    \begin{bmatrix}\beta\\u_o\end{bmatrix}^{\!\top}
    r_o^z
    +c_o,
    \label{eq:joint_old_likelihood_ratio}
\end{equation}
where
\begin{equation}
    R_o^z
    =
    \begin{bmatrix}
        R_{\beta\beta,o} & R_{\beta u,o}\\
        R_{u\beta,o} & R_{uu,o}
    \end{bmatrix},
    \qquad
    r_o^z=
    \begin{bmatrix}
        r_{\beta,o}\\r_{u,o}
    \end{bmatrix}.
    \label{eq:joint_old_natural_blocks}
\end{equation}
The old $u$--$u$ likelihood precision will be maintained in the structured form
\begin{equation}
    R_{uu,o}=B_o\otimes G,
    \label{eq:joint_old_Ruu_kron}
\end{equation}
while the cross block $R_{\beta u,o}$ is retained as a $d\times M_tM_s$ matrix.
Since $d$ is small, this cross block is affordable and is the key object that
preserves $\beta$--$u$ posterior dependence.

Under the GP prior,
\begin{equation}
    p(u_o\mid u_n)
    =
    \mathcal{N}(L_{on}u_n,\Sigma_{o\mid n}),
    \qquad
    L_{on}=K_{on}K_{nn}^{-1}.
    \label{eq:joint_old_given_new}
\end{equation}
Introduce the deterministic transfer matrix for the conditional mean
\begin{equation}
    M_{on}
    :=
    \begin{bmatrix}
        I_d & 0\\
        0 & L_{on}
    \end{bmatrix}.
    \label{eq:joint_transfer_matrix}
\end{equation}
Taking the expectation in \eqref{eq:joint_streaming_vfe} gives, up to constants,
\begin{equation}
    \E_{p(u_o\mid u_n)}
    \left[
        \log
        \frac{q_{n-1}(\beta,u_o)}{p(\beta)p(u_o)}
    \right]
    =
    -\frac{1}{2}
    \begin{bmatrix}\beta\\u_n\end{bmatrix}^{\!\top}
    R_{o\rightarrow n}^z
    \begin{bmatrix}\beta\\u_n\end{bmatrix}
    +
    \begin{bmatrix}\beta\\u_n\end{bmatrix}^{\!\top}
    r_{o\rightarrow n}^z
    +\mathrm{const},
\end{equation}
where
\begin{equation}
    R_{o\rightarrow n}^z=M_{on}^\top R_o^zM_{on},
    \qquad
    r_{o\rightarrow n}^z=M_{on}^\top r_o^z .
    \label{eq:joint_old_transfer_natural_compact}
\end{equation}
The same transfer can be written in block form as
\begin{equation}
    R_{o\rightarrow n}^z
    =
    \begin{bmatrix}
        R_{\beta\beta,o\rightarrow n} & R_{\beta u,o\rightarrow n}\\
        R_{u\beta,o\rightarrow n} & R_{uu,o\rightarrow n}
    \end{bmatrix},
    \qquad
    r_{o\rightarrow n}^z
    =
    \begin{bmatrix}
        r_{\beta,o\rightarrow n}\\
        r_{u,o\rightarrow n}
    \end{bmatrix}.
    \label{eq:joint_old_transfer_blocks_matrix}
\end{equation}
The conditional covariance term
$\tr(R_{uu,o}\Sigma_{o\mid n})$ is absorbed into the constant because it does not
depend on $(\beta,u_n)$.

Expanding the transferred blocks gives
\begin{align}
    R_{\beta\beta,o\rightarrow n}
    &=R_{\beta\beta,o},
    \label{eq:transfer_Rbb}\\
    R_{\beta u,o\rightarrow n}
    &=R_{\beta u,o}L_{on},
    \label{eq:transfer_Rbu}\\
    R_{u\beta,o\rightarrow n}
    &=L_{on}^\top R_{u\beta,o},
    \label{eq:transfer_Rub}\\
    R_{uu,o\rightarrow n}
    &=L_{on}^\top R_{uu,o}L_{on},
    \label{eq:transfer_Ruu}\\
    r_{\beta,o\rightarrow n}
    &=r_{\beta,o},
    \label{eq:transfer_rb}\\
    r_{u,o\rightarrow n}
    &=L_{on}^\top r_{u,o}.
    \label{eq:transfer_ru}
\end{align}
These equations are the route-B replacement for the earlier mean-field transfer: the
$\beta$--$u$ cross natural parameter is explicitly propagated.

\subsection{Assimilating the New Block}
\label{subsec:joint_new_block_assimilation}

For the new block, define
\begin{equation}
    H_n=
    \begin{bmatrix}
        \Phi_n & A_n
    \end{bmatrix}.
\end{equation}
The Gaussian likelihood contributes the joint natural parameters
\begin{equation}
    R_{\mathrm{new},n}^z
    =
    \frac{1}{\sigma^2}H_n^\top H_n,
    \qquad
    r_{\mathrm{new},n}^z
    =
    \frac{1}{\sigma^2}H_n^\top y_n.
\end{equation}
Equivalently, in block-matrix form,
\begin{equation}
    R_{\mathrm{new},n}^z
    =
    \begin{bmatrix}
        R_{\beta\beta,\mathrm{new}} & R_{\beta u,\mathrm{new}}\\
        R_{u\beta,\mathrm{new}} & R_{uu,\mathrm{new}}
    \end{bmatrix},
    \qquad
    r_{\mathrm{new},n}^z
    =
    \begin{bmatrix}
        r_{\beta,\mathrm{new}}\\
        r_{u,\mathrm{new}}
    \end{bmatrix}.
    \label{eq:new_block_joint_blocks_matrix}
\end{equation}
The individual blocks are
\begin{align}
    R_{\beta\beta,\mathrm{new}}
    &=\frac{1}{\sigma^2}\Phi_n^\top\Phi_n,\\
    R_{\beta u,\mathrm{new}}
    &=\frac{1}{\sigma^2}\Phi_n^\top A_n,\\
    R_{u\beta,\mathrm{new}}
    &=\frac{1}{\sigma^2}A_n^\top\Phi_n,\\
    R_{uu,\mathrm{new}}
    &=\frac{1}{\sigma^2}A_n^\top A_n,
\end{align}
and
\begin{align}
    r_{\beta,\mathrm{new}}
    &=\frac{1}{\sigma^2}\Phi_n^\top y_n,\\
    r_{u,\mathrm{new}}
    &=\frac{1}{\sigma^2}A_n^\top y_n.
\end{align}
The accumulated likelihood natural parameters in the new coordinates are therefore
\begin{equation}
    R_n^z
    =
    \begin{bmatrix}
        R_{\beta\beta,n} & R_{\beta u,n}\\
        R_{u\beta,n} & R_{uu,n}
    \end{bmatrix},
    \qquad
    r_n^z
    =
    \begin{bmatrix}
        r_{\beta,n}\\
        r_{u,n}
    \end{bmatrix}.
    \label{eq:joint_accumulated_blocks_matrix}
\end{equation}
The block-wise accumulated likelihood statistics are
\begin{align}
    R_{\beta\beta,n}
    &=R_{\beta\beta,o\rightarrow n}+\frac{1}{\sigma^2}\Phi_n^\top\Phi_n
      =R_{\beta\beta,o}+\frac{1}{\sigma^2}\Phi_n^\top\Phi_n,
    \label{eq:joint_accum_Rbb}\\
    R_{\beta u,n}
    &=R_{\beta u,o\rightarrow n}+\frac{1}{\sigma^2}\Phi_n^\top A_n
      =R_{\beta u,o}L_{on}+\frac{1}{\sigma^2}\Phi_n^\top A_n,
    \label{eq:joint_accum_Rbu}\\
    R_{u\beta,n}
    &=R_{u\beta,o\rightarrow n}+\frac{1}{\sigma^2}A_n^\top\Phi_n
      =L_{on}^\top R_{u\beta,o}+\frac{1}{\sigma^2}A_n^\top\Phi_n,
    \label{eq:joint_accum_Rub}\\
    R_{uu,n}
    &=R_{uu,o\rightarrow n}+\frac{1}{\sigma^2}A_n^\top A_n
      =L_{on}^\top R_{uu,o}L_{on}+\frac{1}{\sigma^2}A_n^\top A_n,
    \label{eq:joint_accum_Ruu}\\
    r_{\beta,n}
    &=r_{\beta,o\rightarrow n}+\frac{1}{\sigma^2}\Phi_n^\top y_n
      =r_{\beta,o}+\frac{1}{\sigma^2}\Phi_n^\top y_n,
    \label{eq:joint_accum_rb}\\
    r_{u,n}
    &=r_{u,o\rightarrow n}+\frac{1}{\sigma^2}A_n^\top y_n
      =L_{on}^\top r_{u,o}+\frac{1}{\sigma^2}A_n^\top y_n.
    \label{eq:joint_accum_ru}
\end{align}

Adding the base priors gives the posterior information parameters
\begin{align}
    \Lambda_{\beta\beta,n}
    &=P_\beta+R_{\beta\beta,n},
    \label{eq:joint_Lbb_total}\\
    \Lambda_{\beta u,n}
    &=R_{\beta u,n},
    \label{eq:joint_Lbu_total}\\
    \Lambda_{uu,n}
    &=K_{nn}^{-1}+R_{uu,n},
    \label{eq:joint_Luu_total}\\
    h_{\beta,n}
    &=P_\beta m_{\beta,0}+r_{\beta,n},
    \label{eq:joint_hb_total}\\
    h_{u,n}
    &=r_{u,n}.
    \label{eq:joint_hu_total}
\end{align}
This is the structured joint natural-parameter update used by the main method.
A slowly drifting $\beta$ can be introduced as an approximation by replacing
$(m_{\beta,0},K_\beta)$ with a predicted Gaussian prior at each block. The exact
old-likelihood-ratio derivation above is cleanest for a static $\beta$ with a fixed
base prior.

\subsection{Projected-Prior Baseline with Cross-Covariance}
\label{subsec:projected_prior_approximation}

A cheaper ablation is to project the previous joint posterior moments instead of
propagating the old likelihood ratio. If
$q_{n-1}(\beta,u_o)$ has moments
$(m_{\beta,o},m_{u,o},S_{\beta\beta,o},S_{\beta u,o},S_{uu,o})$, then
\begin{equation}
    p_n^{\mathrm{proj}}(\beta,u_n)
    =
    \int p(u_n\mid u_o)q_{n-1}(\beta,u_o)\,du_o
\end{equation}
with $L_{no}:=K_{no}K_{oo}^{-1}$ gives
\begin{align}
    m_{\beta,n\mid n-1}^{\mathrm{proj}}&=m_{\beta,o},\\
    m_{u,n\mid n-1}^{\mathrm{proj}}&=L_{no}m_{u,o},\\
    S_{\beta\beta,n\mid n-1}^{\mathrm{proj}}&=S_{\beta\beta,o},\\
    S_{\beta u,n\mid n-1}^{\mathrm{proj}}&=S_{\beta u,o}L_{no}^\top,\\
    S_{uu,n\mid n-1}^{\mathrm{proj}}
    &=K_{nn}+L_{no}(S_{uu,o}-K_{oo})L_{no}^\top.
\end{align}
This projected-prior update is not the main theoretical method, but it is useful as a
baseline because it tests whether the old-likelihood-ratio transfer is better than a
simple filtering-style moment projection.

\subsection{Summary of the Structured Changing-Basis Transfer}
\label{subsec:changing_basis_summary}

The changing-basis update now transfers the old joint likelihood statistic rather than
only the marginal inducing statistic. The essential update is
\begin{equation}
    \begin{bmatrix}
        R_{\beta\beta,o\rightarrow n} & R_{\beta u,o\rightarrow n}\\
        R_{u\beta,o\rightarrow n} & R_{uu,o\rightarrow n}
    \end{bmatrix}
    =
    \begin{bmatrix}
        I & 0\\0 & L_{on}
    \end{bmatrix}^{\!\top}
    \begin{bmatrix}
        R_{\beta\beta,o} & R_{\beta u,o}\\
        R_{u\beta,o} & R_{uu,o}
    \end{bmatrix}
    \begin{bmatrix}
        I & 0\\0 & L_{on}
    \end{bmatrix}.
\end{equation}
The $u$--$u$ block remains Kronecker/Sylvester structured under the conditions derived
in Section~\ref{sec:kronecker_aware_implementation}, while the cross block
$R_{\beta u}$ is retained explicitly. This is the route-B main-paper version: it fixes
stale residual targets and keeps the structured $\beta$--$u$ covariance needed for
well-calibrated predictive uncertainty.

\section{Structured Joint Posterior Recovery by Schur Complement}
\label{sec:structured_joint_schur}

After the natural-parameter update, we need posterior means and selected covariance
quadratic forms without inverting the full joint precision. Write
\begin{equation}
    \Lambda_n=
    \begin{bmatrix}
        A_\beta & B_{\beta u}\\
        B_{\beta u}^\top & D_u
    \end{bmatrix},
    \qquad
    h_n=
    \begin{bmatrix}
        h_\beta\\h_u
    \end{bmatrix},
    \label{eq:schur_block_notation}
\end{equation}
where
\begin{equation}
    A_\beta=\Lambda_{\beta\beta,n},
    \qquad
    B_{\beta u}=\Lambda_{\beta u,n},
    \qquad
    D_u=\Lambda_{uu,n}.
\end{equation}
The matrix $D_u$ is large but has the Kronecker/Sylvester form. The dimension of
$A_\beta$ is only $d\times d$.

Define the Schur-complement precision
\begin{equation}
    \Lambda_{\beta\mid u}
    :=
    A_\beta-B_{\beta u}D_u^{-1}B_{\beta u}^\top.
    \label{eq:beta_schur_complement}
\end{equation}
This object is a precision matrix for the low-dimensional $\beta$ block after
eliminating the GP block; it is not the posterior covariance of $\beta$.
To form it, solve
\begin{equation}
    D_u W=B_{\beta u}^\top,
    \qquad
    D_uv_h=h_u,
    \label{eq:schur_sylvester_solves}
\end{equation}
using the Sylvester solver. Since $B_{\beta u}^\top$ has only $d$ columns, this requires
only $d+1$ Sylvester solves per block. Then
\begin{equation}
    \Lambda_{\beta\mid u}=A_\beta-B_{\beta u}W.
\end{equation}
The posterior mean of $\beta$ is
\begin{equation}
    m_\beta
    =
    \Lambda_{\beta\mid u}^{-1}
    \left(
        h_\beta-B_{\beta u}v_h
    \right),
    \label{eq:schur_beta_mean}
\end{equation}
and the posterior mean of $u$ is
\begin{equation}
    m_u
    =
    v_h-Wm_\beta.
    \label{eq:schur_u_mean}
\end{equation}
The covariance blocks needed for prediction are not materialized densely. They are
characterized by
\begin{align}
    S_{\beta\beta}&=\Lambda_{\beta\mid u}^{-1},
    \label{eq:Sbb_schur}\\
    S_{\beta u}&=-\Lambda_{\beta\mid u}^{-1}B_{\beta u}D_u^{-1},
    \label{eq:Sbu_schur}\\
    S_{uu}&=D_u^{-1}+D_u^{-1}B_{\beta u}^\top
    \Lambda_{\beta\mid u}^{-1}B_{\beta u}D_u^{-1}.
    \label{eq:Suu_schur}
\end{align}
Thus the route-B method keeps the $\beta$--$u$ posterior dependence while preserving the
large-scale solver for the GP block.


\section{Kronecker-Aware Implementation}
\label{sec:kronecker_aware_implementation}

The Kronecker form
\begin{equation}
    A_n=T_n\otimes C
\end{equation}
should be used explicitly. Therefore,
\begin{equation}
    A_n^\top A_n
    =
    (T_n^\top T_n)\otimes(C^\top C).
\end{equation}
Let
\begin{equation}
    G:=C^\top C.
\end{equation}
Since the spatial inducing set is fixed, $G$ can be precomputed.

\subsection{Kronecker-Preserving $u$--$u$ Old-Likelihood Transfer}
\label{subsec:kronecker_preserving_transfer}

Let
\begin{equation}
    K_{nn}=K_{nn}^{(t)}\otimes K_s,
    \qquad
    K_{oo}=K_{oo}^{(t)}\otimes K_s,
    \qquad
    K_{on}=K_{on}^{(t)}\otimes K_s,
\end{equation}
where $K_s=K_{ZZ}^{(s)}$. Since the spatial inducing locations are fixed,
\begin{equation}
    L_{on}=K_{on}K_{nn}^{-1}=L_{on}^{(t)}\otimes I_s,
    \qquad
    L_{on}^{(t)}:=K_{on}^{(t)}(K_{nn}^{(t)})^{-1}.
    \label{eq:temporal_only_transfer_operator}
\end{equation}
If the old $u$--$u$ likelihood precision is maintained as
\begin{equation}
    R_{uu,o}=B_o\otimes G,
\end{equation}
then
\begin{align}
    R_{uu,o\rightarrow n}
    &=L_{on}^\top R_{uu,o}L_{on} \\
    &=
    \left[(L_{on}^{(t)})^\top B_oL_{on}^{(t)}\right]\otimes G.
\end{align}
Define
\begin{equation}
    B_{o\rightarrow n}:=(L_{on}^{(t)})^\top B_oL_{on}^{(t)}.
\end{equation}
The new block contributes
\begin{equation}
    \frac{1}{\sigma^2}A_n^\top A_n
    =
    \frac{1}{\sigma^2}(T_n^\top T_n)\otimes G.
\end{equation}
Therefore the updated temporal likelihood statistic is
\begin{equation}
    B_n
    =
    B_{o\rightarrow n}
    +
    \frac{1}{\sigma^2}T_n^\top T_n.
    \label{eq:temporal_stat_ssgp_transfer}
\end{equation}
The $u$--$u$ precision block in the full joint posterior is
\begin{equation}
    \Lambda_{uu,n}
    =
    (K_{nn}^{(t)})^{-1}\otimes K_s^{-1}
    +
    B_n\otimes G.
    \label{eq:kronecker_precision_after_ssgp_transfer}
\end{equation}
This is the same Sylvester-compatible structure as in the marginal GP route, but now it
is the lower-right block of a structured joint precision over $(\beta,u)$.

\subsection{Structured $\beta$--$u$ Cross Block}
\label{subsec:structured_beta_u_cross_block}

Route B additionally maintains the cross natural parameter
\begin{equation}
    R_{\beta u,n}\in\mathbb{R}^{d\times M_tM_s}.
\end{equation}
The old cross block transfers as
\begin{equation}
    R_{\beta u,o\rightarrow n}
    =
    R_{\beta u,o}L_{on}
    =
    R_{\beta u,o}(L_{on}^{(t)}\otimes I_s).
    \label{eq:cross_block_transfer_kron}
\end{equation}
This should be implemented without materializing $L_{on}$ when possible: each row of
$R_{\beta u,o}$ can be reshaped into an $M_s\times M_t$ matrix and multiplied along the
temporal dimension by $L_{on}^{(t)}$ according to the vectorization convention.

The new likelihood adds
\begin{equation}
    R_{\beta u,\mathrm{new}}
    =
    \frac{1}{\sigma^2}\Phi_n^\top A_n
    =
    \frac{1}{\sigma^2}\Phi_n^\top(T_n\otimes C).
    \label{eq:new_cross_block_likelihood}
\end{equation}
The updated cross block is
\begin{equation}
    R_{\beta u,n}
    =
    R_{\beta u,o}(L_{on}^{(t)}\otimes I_s)
    +
    \frac{1}{\sigma^2}\Phi_n^\top A_n.
    \label{eq:updated_cross_block}
\end{equation}
This block is a likelihood natural-precision cross block, not the posterior
cross-covariance itself. The posterior cross-covariance is recovered implicitly as
\begin{equation}
    S_{\beta u}
    =
    -\Lambda_{\beta\mid u}^{-1}B_{\beta u}D_u^{-1}.
    \label{eq:posterior_cross_covariance_implicit}
\end{equation}
The cross block need not be a single Kronecker product. That is acceptable because its row
dimension is only $d$. It affects posterior recovery only through matrix products with
$D_u^{-1}$, which are computed by repeated Sylvester solves.

The corresponding information vector transfer is
\begin{equation}
    r_{u,o\rightarrow n}=L_{on}^\top r_{u,o},
    \qquad
    r_{\beta,o\rightarrow n}=r_{\beta,o},
\end{equation}
and the new information contributions are
\begin{equation}
    r_{u,\mathrm{new}}=\frac{1}{\sigma^2}A_n^\top y_n,
    \qquad
    r_{\beta,\mathrm{new}}=\frac{1}{\sigma^2}\Phi_n^\top y_n.
\end{equation}

\subsection{Sylvester Solves for the $u$ Block}
\label{subsec:sylvester_reformulation}

With the Kronecker-preserving transfer, linear systems involving
$D_u=\Lambda_{uu,n}$ reduce to Sylvester equations. Let
\begin{equation}
    q=\vecop(Q),
    \qquad
    z=D_u^{-1}q=\vecop(Z),
\end{equation}
where $Q,Z\in\mathbb{R}^{M_s\times M_t}$. Using
\begin{equation}
    D_u=(K_{nn}^{(t)})^{-1}\otimes K_s^{-1}+B_n\otimes G,
\end{equation}
and the identity $(A\otimes B)\vecop(Z)=\vecop(BZA^\top)$, the system
$D_uz=q$ is equivalent to
\begin{equation}
    K_s^{-1}Z(K_{nn}^{(t)})^{-1}+GZB_n=Q,
    \label{eq:sylvester_after_ssgp_transfer}
\end{equation}
assuming symmetric temporal matrices. The same solver is used for
\begin{equation}
    D_u^{-1}h_u,
    \qquad
    D_u^{-1}R_{u\beta,n},
    \qquad
    D_u^{-1}a_*.
\end{equation}
The second solve has only $d$ right-hand sides and is what makes the Schur complement
for $\beta$ computationally practical.

\begin{remark}
If the spatial observation pattern or spatial projection changes across online blocks,
the $u$--$u$ precision may become a sum of Kronecker products rather than a single
Sylvester form. The structured joint route still applies, but $D_u^{-1}$ must then be
implemented with a generalized Sylvester or conjugate-gradient solver using Kronecker
matrix-vector products.
\end{remark}


\section{Prediction and Uncertainty}
\label{sec:prediction_uncertainty}

For a test point $x_*=(t_*,s_*)$, let
\begin{equation}
    \phi_*:=\phi(t_*,s_*),
    \qquad
    a_*^\top:=K_{*u}K_{uu}^{-1}.
\end{equation}
Let the structured joint precision be partitioned as in
\eqref{eq:schur_block_notation}, with $D_u=\Lambda_{uu}$ and
$B_{\beta u}=\Lambda_{\beta u}$. The predictive mean is
\begin{equation}
    \E[y_*]
    =
    \phi_*^\top m_\beta+a_*^\top m_u.
    \label{eq:structured_predictive_mean}
\end{equation}

The sparse GP conditional variance term is
\begin{equation}
    \nu_*
    :=
    k(x_*,x_*)-K_{*u}K_{uu}^{-1}K_{u*}.
\end{equation}
The remaining posterior-parameter uncertainty is
\begin{equation}
    \begin{bmatrix}\phi_*\\a_*\end{bmatrix}^{\!\top}
    \Lambda^{-1}
    \begin{bmatrix}\phi_*\\a_*\end{bmatrix}.
\end{equation}
This quantity can be evaluated without materializing the full covariance. Solve
\begin{equation}
    D_uv_*=a_*.
\end{equation}
Then, using the Schur-complement precision
$\Lambda_{\beta\mid u}=A_\beta-B_{\beta u}D_u^{-1}B_{\beta u}^\top$, we have
\begin{equation}
    \begin{bmatrix}\phi_*\\a_*\end{bmatrix}^{\!\top}
    \Lambda^{-1}
    \begin{bmatrix}\phi_*\\a_*\end{bmatrix}
    =
    a_*^\top v_*
    +
    \left(\phi_*-B_{\beta u}v_*\right)^\top
    \Lambda_{\beta\mid u}^{-1}
    \left(\phi_*-B_{\beta u}v_*\right).
    \label{eq:structured_predictive_quadratic}
\end{equation}
Therefore the predictive variance is
\begin{equation}
    \Var(y_*)
    =
    \sigma^2+
    \nu_*+
    a_*^\top v_*
    +
    \left(\phi_*-B_{\beta u}v_*\right)^\top
    \Lambda_{\beta\mid u}^{-1}
    \left(\phi_*-B_{\beta u}v_*\right).
    \label{eq:structured_predictive_variance}
\end{equation}
Equation~\eqref{eq:structured_predictive_variance} contains the effect of the
$\beta$--$u$ cross covariance implicitly. Expanding the full covariance would give
\begin{equation}
    \phi_*^\top S_{\beta\beta}\phi_*
    +a_*^\top S_{uu}a_*
    +2\phi_*^\top S_{\beta u}a_*.
\end{equation}
The cross term $2\phi_*^\top S_{\beta u}a_*$ is absent under mean-field. Route B keeps
it through the Schur-complement formula above while still requiring only Sylvester solves
with the $u$ block.

For comparison, the mean-field ablation $q(\beta)q(u)$ uses
\begin{equation}
    \Var_{\mathrm{MF}}(y_*)
    =
    \sigma^2+
    \nu_*+
    \phi_*^\top S_{\beta,\mathrm{MF}}\phi_*
    +
    a_*^\top S_{u,\mathrm{MF}}a_* ,
\end{equation}
which drops the cross-covariance contribution and should therefore be treated as a
calibration ablation rather than the main predictive distribution.


\section{Experiments and Ablations}
\label{sec:experiments_ablations}

The recommended experimental sequence should isolate the contribution of each modeling
choice: joint training, changing-basis transfer, old-likelihood-ratio transfer,
Kronecker structure, and retention of structured $\beta$--$u$ covariance.

\subsection{Methods to Compare}

At minimum, compare the following methods:
\begin{enumerate}[label=(\arabic*)]
    \item \textbf{Two-stage future-only baseline.}
    The linear mean is updated first and the GP is trained on current residuals only.

    \item \textbf{Two-stage residual GP with buffer or sketch correction.}
    This tests whether explicit residual repair can mitigate stale residual targets.

    \item \textbf{Fixed-basis exact joint Gaussian baseline.}
    This uses a fixed inducing representation and performs the exact joint Gaussian
    update over $(\beta,u)$.

    \item \textbf{Joint mean-field with projected-prior transfer.}
    This is the efficient filtering-style approximation but drops $\beta$--$u$ posterior
    covariance.

    \item \textbf{Joint mean-field with SSGP-style old-likelihood-ratio transfer.}
    This tests whether old-likelihood transfer helps even without the structured cross
    covariance.

    \item \textbf{Structured joint posterior with SSGP-style transfer.}
    This is the main route-B method. It keeps $\Lambda_{\beta u}$ and recovers posterior
    moments using the Schur complement.

    \item \textbf{Structured joint posterior without changing-basis transfer.}
    This ablation keeps $\Lambda_{\beta u}$ but ignores that $u_o$ and $u_n$ live in
    different HiPPO coordinates.
\end{enumerate}

\subsection{Ablations}

Important ablations include:
\begin{enumerate}[label=(\arabic*)]
    \item number of temporal HiPPO basis functions $M_t$;
    \item number of spatial inducing points $M_s$;
    \item number of RFF frequencies $R$;
    \item fixed temporal basis versus changing HiPPO basis;
    \item no transfer versus projected-prior transfer versus old-likelihood-ratio transfer;
    \item maintaining $R_{uu,o}$ as dense versus maintaining $B_o\otimes G$;
    \item dropping versus retaining the cross block $\Lambda_{\beta u}$;
    \item exact fixed-basis joint posterior versus structured Schur approximation;
    \item static $\beta$ versus approximate drifting $\beta$;
    \item PSD projection, damping, or jitter for old-likelihood natural precisions when
    dense approximations are used.
\end{enumerate}

\subsection{Metrics}

The evaluation should not rely only on RMSE. Report:
\begin{enumerate}[label=(\arabic*)]
    \item RMSE;
    \item MAE;
    \item test negative log likelihood;
    \item predictive interval coverage, for example $50\%$, $90\%$, and $95\%$;
    \item runtime per online block;
    \item peak memory;
    \item RMSE/NLL under matched memory budgets;
    \item calibration error for Gaussian predictive intervals;
    \item error on continual evaluation windows, not only future-only forecasting.
\end{enumerate}

The most important comparison is whether the structured route-B method improves
continual evaluation and uncertainty calibration relative to residual-based learning,
projected-prior transfer, and mean-field SSGP-style transfer.


\section{Summary}
\label{sec:summary}

The proposed model follows the logical structure
\begin{equation}
\begin{aligned}
    &\text{two-stage learning has stale residuals} \\
    \Rightarrow\;&
    \text{joint training fixes the observation target} \\
    \Rightarrow\;&
    \text{changing HiPPO bases require old-information transfer} \\
    \Rightarrow\;&
    \text{SSGP-style old-likelihood-ratio transfer gives a principled streaming update} \\
    \Rightarrow\;&
    \text{route B keeps the structured }\beta\text{--}u\text{ covariance} \\
    \Rightarrow\;&
    \text{Kronecker/Sylvester structure keeps the method scalable.}
\end{aligned}
\end{equation}

The observation model is
\begin{equation}
    y_n=\Phi_n\beta+A_nu_n+\epsilon_n,
    \qquad
    \epsilon_n\sim\mathcal{N}(0,\sigma^2I).
\end{equation}
Unlike residual-based learning, the likelihood is always applied to $y_n$:
\begin{equation}
    p(y_n\mid \beta,u_n)
    =
    \mathcal{N}(y_n\mid \Phi_n\beta+A_nu_n,\sigma^2I).
\end{equation}

For a changing HiPPO temporal basis, the main update propagates the old joint likelihood
ratio
\begin{equation}
    \frac{q_{n-1}(\beta,u_o)}{p(\beta)p(u_o)}.
\end{equation}
If the old likelihood natural parameters are
\begin{equation}
    R_o^z=
    \begin{bmatrix}
        R_{\beta\beta,o} & R_{\beta u,o}\\
        R_{u\beta,o} & R_{uu,o}
    \end{bmatrix},
    \qquad
    r_o^z=
    \begin{bmatrix}
        r_{\beta,o}\\r_{u,o}
    \end{bmatrix},
\end{equation}
and $p(u_o\mid u_n)=\mathcal{N}(L_{on}u_n,\Sigma_{o\mid n})$, then the old likelihood is
transferred by
\begin{equation}
    R_{o\rightarrow n}^z
    =
    \begin{bmatrix}
        I & 0\\0 & L_{on}
    \end{bmatrix}^{\!\top}
    R_o^z
    \begin{bmatrix}
        I & 0\\0 & L_{on}
    \end{bmatrix},
    \qquad
    r_{o\rightarrow n}^z
    =
    \begin{bmatrix}
        I & 0\\0 & L_{on}
    \end{bmatrix}^{\!\top}
    r_o^z.
\end{equation}
The resulting accumulated likelihood blocks are
\begin{align}
    R_{\beta\beta,n}
    &=R_{\beta\beta,o}+\frac{1}{\sigma^2}\Phi_n^\top\Phi_n,\\
    R_{\beta u,n}
    &=R_{\beta u,o}L_{on}+\frac{1}{\sigma^2}\Phi_n^\top A_n,\\
    R_{uu,n}
    &=L_{on}^\top R_{uu,o}L_{on}+\frac{1}{\sigma^2}A_n^\top A_n.
\end{align}
After adding the priors, the posterior precision has the route-B structure
\begin{equation}
    \Lambda_n=
    \begin{bmatrix}
        \Lambda_{\beta\beta,n} & \Lambda_{\beta u,n}\\
        \Lambda_{u\beta,n} & \Lambda_{uu,n}
    \end{bmatrix}.
\end{equation}
The $u$--$u$ block keeps the Sylvester form. Since the changing basis acts only in the
temporal dimension and the spatial inducing set is fixed,
\begin{equation}
    L_{on}=L_{on}^{(t)}\otimes I_s.
\end{equation}
If
\begin{equation}
    R_{uu,o}=B_o\otimes G,
    \qquad
    G=C^\top C,
\end{equation}
then
\begin{equation}
    L_{on}^\top R_{uu,o}L_{on}
    =
    \left[(L_{on}^{(t)})^\top B_oL_{on}^{(t)}\right]\otimes G.
\end{equation}
Thus
\begin{equation}
    \Lambda_{uu,n}
    =
    (K_{nn}^{(t)})^{-1}\otimes K_s^{-1}+B_n\otimes G,
\end{equation}
where
\begin{equation}
    B_n=(L_{on}^{(t)})^\top B_oL_{on}^{(t)}+\frac{1}{\sigma^2}T_n^\top T_n.
\end{equation}
The Sylvester equation
\begin{equation}
    K_s^{-1}Z(K_{nn}^{(t)})^{-1}+GZB_n=Q
\end{equation}
therefore remains available for $D_u^{-1}$ solves.

Posterior moments are recovered by Schur complement. For
\begin{equation}
    \Lambda_n=
    \begin{bmatrix}
        A_\beta & B_{\beta u}\\
        B_{\beta u}^\top & D_u
    \end{bmatrix},
\end{equation}
define the Schur-complement precision
\begin{equation}
    \Lambda_{\beta\mid u}=A_\beta-B_{\beta u}D_u^{-1}B_{\beta u}^\top.
\end{equation}
Then
\begin{equation}
    m_\beta=\Lambda_{\beta\mid u}^{-1}(h_\beta-B_{\beta u}D_u^{-1}h_u),
    \qquad
    m_u=D_u^{-1}(h_u-B_{\beta u}^\top m_\beta).
\end{equation}
Only $d+1$ Sylvester solves are needed to form the Schur complement and posterior means.
For prediction, the structured posterior variance is evaluated as
\begin{equation}
    \Var(y_*)
    =
    \sigma^2+\nu_*+a_*^\top v_*
    +
    (\phi_*-B_{\beta u}v_*)^\top \Lambda_{\beta\mid u}^{-1}(\phi_*-B_{\beta u}v_*),
    \qquad
    D_uv_*=a_*.
\end{equation}
This formula retains the $\beta$--$u$ covariance term that the mean-field approximation
would drop. Therefore the main theoretical contribution is no longer only
old-likelihood-ratio transfer for $u$, but \emph{structured joint old-likelihood transfer
with a Kronecker-preserving GP block and a Schur-complement correction for the linear
mean}.

\appendix

\section{Proofs for Kronecker-Preserving Changing-Basis Transfer}
\label{app:kron_preserving_transfer_proofs}

This appendix gives the detailed proof of the two structural claims used
in the main text:
\begin{equation}
    L_{on}=L_{on}^{(t)}\otimes I_s,
    \label{eq:app_Lon_kron_claim}
\end{equation}
and
\begin{equation}
    R_{uu,o}=B_o\otimes G.
    \label{eq:app_Ruu_kron_claim}
\end{equation}
The first identity is a consequence of the separable spatio-temporal
kernel and the fact that the spatial inducing set is fixed. The second
identity is not true for an arbitrary dense posterior covariance; it is an
algorithmic invariant obtained by maintaining the old likelihood natural
precision in a Kronecker-structured form.

\subsection{\texorpdfstring{Proof that \(L_{on}=L_{on}^{(t)}\otimes I_s\)}{Proof that Lon = Lon(t) tensor Is}}
\label{app:proof_Lon_temporal_only}

Assume the spatio-temporal kernel is separable:
\begin{equation}
    k\big((t,s),(t',s')\big)
    =
    k_t(t,t')k_s(s,s').
    \label{eq:app_sep_kernel}
\end{equation}
Let the mixed HiPPO interdomain inducing variables be
\begin{equation}
    u_{\ell,j}(T)
    =
    \int_0^T
    g_\ell^{(T)}(x)f(x,z_j^{(s)})\,dx,
    \label{eq:app_mixed_inducing}
\end{equation}
where \(\ell\) indexes the temporal HiPPO basis and \(j\) indexes the
fixed spatial inducing location \(z_j^{(s)}\).

Let
\begin{equation}
    u_o := u(T_o),
    \qquad
    u_n := u(T_n),
\end{equation}
where \(T_o=T_{n-1}\) is the old temporal horizon and \(T_n\) is the new
temporal horizon. The temporal HiPPO basis changes from
\(g_\ell^{(T_o)}\) to \(g_\ell^{(T_n)}\), but the spatial inducing
locations \(Z_s=\{z_1^{(s)},\ldots,z_{M_s}^{(s)}\}\) are fixed.

For two inducing entries,
\begin{equation}
    u_{\ell,j}^{old}
    =
    \int_0^{T_o}
    g_\ell^{(T_o)}(x)f(x,z_j^{(s)})\,dx,
\end{equation}
and
\begin{equation}
    u_{m,j'}^{new}
    =
    \int_0^{T_n}
    g_m^{(T_n)}(x')f(x',z_{j'}^{(s)})\,dx',
\end{equation}
their prior covariance is
\begin{align}
    \cov
    \left(
        u_{\ell,j}^{old},
        u_{m,j'}^{new}
    \right)
    &=
    \int_0^{T_o}
    \int_0^{T_n}
    g_\ell^{(T_o)}(x)
    g_m^{(T_n)}(x')
    k\big((x,z_j^{(s)}),(x',z_{j'}^{(s)})\big)
    \,dx'\,dx                                                  \\
    &=
    \int_0^{T_o}
    \int_0^{T_n}
    g_\ell^{(T_o)}(x)
    g_m^{(T_n)}(x')
    k_t(x,x')
    k_s(z_j^{(s)},z_{j'}^{(s)})
    \,dx'\,dx                                                   \\
    &=
    \left[
    \int_0^{T_o}
    \int_0^{T_n}
    g_\ell^{(T_o)}(x)
    g_m^{(T_n)}(x')
    k_t(x,x')
    \,dx'\,dx
    \right]
    k_s(z_j^{(s)},z_{j'}^{(s)}).
    \label{eq:app_old_new_cov_entry}
\end{align}
Thus the old-new covariance factorizes into a temporal covariance factor
and a spatial covariance factor. With the vectorization convention used in
the main text,
\begin{equation}
    K_{on}
    =
    K_{on}^{(t)}\otimes K_s,
    \qquad
    K_s:=K_{ZZ}^{(s)}.
    \label{eq:app_Kon_kron}
\end{equation}
Similarly,
\begin{equation}
    K_{nn}
    =
    K_{nn}^{(t)}\otimes K_s,
    \qquad
    K_{oo}
    =
    K_{oo}^{(t)}\otimes K_s.
    \label{eq:app_Knn_Koo_kron}
\end{equation}

The conditional mean operator in
\begin{equation}
    p(u_o\mid u_n)
    =
    \mathcal{N}
    \left(
        L_{on}u_n,
        \Sigma_{o\mid n}
    \right)
\end{equation}
is
\begin{equation}
    L_{on}
    =
    K_{on}K_{nn}^{-1}.
\end{equation}
Using \eqref{eq:app_Kon_kron}--\eqref{eq:app_Knn_Koo_kron} and the
Kronecker identities
\begin{equation}
    (A\otimes B)^{-1}=A^{-1}\otimes B^{-1},
    \qquad
    (A\otimes B)(C\otimes D)=(AC)\otimes(BD),
\end{equation}
we obtain
\begin{align}
    L_{on}
    &=
    \left(
        K_{on}^{(t)}\otimes K_s
    \right)
    \left[
        \left(K_{nn}^{(t)}\right)^{-1}
        \otimes
        K_s^{-1}
    \right]                                                     \\
    &=
    K_{on}^{(t)}
    \left(K_{nn}^{(t)}\right)^{-1}
    \otimes
    K_sK_s^{-1}                                                  \\
    &=
    L_{on}^{(t)}\otimes I_s,
    \label{eq:app_Lon_kron_proof}
\end{align}
where
\begin{equation}
    L_{on}^{(t)}
    :=
    K_{on}^{(t)}
    \left(K_{nn}^{(t)}\right)^{-1}.
\end{equation}
Therefore, under a separable kernel, fixed spatial inducing locations, and
temporal-only changing basis, the old-to-new conditional operator acts
only on the temporal inducing coordinates and leaves the spatial inducing
coordinates unchanged.

\paragraph{Vectorization convention.}
The identity \eqref{eq:app_Lon_kron_proof} assumes the same vectorization
ordering as the Kronecker convention \(K_{uu}=K_{uu}^{(t)}\otimes K_s\).
If an implementation uses the opposite vectorization order, the equivalent
identity becomes \(L_{on}=I_s\otimes L_{on}^{(t)}\). The mathematics is
unchanged, but the code must be consistent with the chosen ordering.

\subsection{\texorpdfstring{Proof and Status of \(R_{uu,o}=B_o\otimes G\)}{Proof and Status of Ruu,o = Bo tensor G}}
\label{app:proof_Ro_kron_natural_stat}

The identity
\begin{equation}
    R_{uu,o}=B_o\otimes G
\end{equation}
should be interpreted carefully. It is not a property of an arbitrary
dense posterior covariance. It is a structured natural-parameter
representation of the accumulated old likelihood information.

For a Gaussian observation model, the likelihood contribution of a block
\(j\) to the inducing-variable precision is
\begin{equation}
    \Lambda_{\mathrm{like},j}
    =
    \frac{1}{\sigma^2}A_j^\top A_j.
    \label{eq:app_likelihood_precision}
\end{equation}
Under the Kronecker HiPPO-SVGP representation,
\begin{equation}
    A_j=T_j\otimes C,
    \label{eq:app_Aj_kron}
\end{equation}
where \(T_j\) is the temporal projection matrix for block \(j\), and
\begin{equation}
    C=K_{XZ}^{(s)}\left(K_{ZZ}^{(s)}\right)^{-1}
\end{equation}
is the fixed spatial projection matrix. Therefore,
\begin{align}
    \Lambda_{\mathrm{like},j}
    &=
    \frac{1}{\sigma^2}
    (T_j\otimes C)^\top(T_j\otimes C)                         \\
    &=
    \frac{1}{\sigma^2}
    (T_j^\top T_j)\otimes(C^\top C).
    \label{eq:app_likelihood_precision_kron}
\end{align}
Define
\begin{equation}
    G:=C^\top C.
\end{equation}
Then each block contributes
\begin{equation}
    \Lambda_{\mathrm{like},j}
    =
    \frac{1}{\sigma^2}
    (T_j^\top T_j)\otimes G.
\end{equation}

The old $u$--$u$ likelihood natural precision before processing block \(n\)
is the sum of historical likelihood precision contributions, expressed in the old
inducing coordinates:
\begin{equation}
    R_{uu,o}
    :=
    \sum_{j<n}
    \Lambda_{\mathrm{like},j}.
\end{equation}
If the spatial projection matrix \(C\) is fixed across the online blocks,
then the spatial factor \(G=C^\top C\) is common to all terms. Hence
\begin{align}
    R_{uu,o}
    &=
    \sum_{j<n}
    \frac{1}{\sigma^2}
    (T_j^\top T_j)\otimes G                                  \\
    &=
    \left[
        \sum_{j<n}
        \frac{1}{\sigma^2}
        T_j^\top T_j
    \right]\otimes G.
\end{align}
Define the temporal old-likelihood precision statistic
\begin{equation}
    B_o
    :=
    \sum_{j<n}
    \frac{1}{\sigma^2}
    T_j^\top T_j.
    \label{eq:app_Bo_definition}
\end{equation}
Then
\begin{equation}
    R_{uu,o}=B_o\otimes G.
    \label{eq:app_Ruu_Bo_G}
\end{equation}

Thus \(R_{uu,o}=B_o\otimes G\) is valid when the algorithm maintains the old
$u$--$u$ likelihood precision directly in natural-parameter form. In a
marginal GP-only Gaussian update, this agrees with
\begin{equation}
    R_{uu,o}=D_{u,o}-K_{oo}^{-1},
\end{equation}
where \(D_{u,o}\) is the lower-right inducing-variable precision block.
In the structured joint route, this block should be maintained as a
natural statistic; it should not be reconstructed from the marginal
covariance of \(u_o\). If one stores an unrestricted dense posterior
covariance, a marginal covariance from a full joint posterior over
\((\beta,u_o)\), or an arbitrary structured approximation, then the dense
quantity obtained by inverting that marginal covariance need not be a
single Kronecker product.

\subsection{Kronecker Preservation of the Old-Likelihood Transfer}
\label{app:proof_Lambda_transfer_kron}

Combining the previous two results proves the Kronecker preservation of
the SSGP-style old-likelihood-ratio transfer.

The transferred old $u$--$u$ likelihood precision is
\begin{equation}
    R_{uu,o\rightarrow n}
    =
    L_{on}^\top R_{uu,o}L_{on}.
\end{equation}
Using
\begin{equation}
    L_{on}=L_{on}^{(t)}\otimes I_s,
    \qquad
    R_{uu,o}=B_o\otimes G,
\end{equation}
we have
\begin{align}
    R_{uu,o\rightarrow n}
    &=
    \left[
        \left(L_{on}^{(t)}\right)^\top\otimes I_s
    \right]
    (B_o\otimes G)
    \left[
        L_{on}^{(t)}\otimes I_s
    \right]                                                     \\
    &=
    \left[
        \left(L_{on}^{(t)}\right)^\top
        B_o
        L_{on}^{(t)}
    \right]
    \otimes G.
\end{align}
Define
\begin{equation}
    B_{o\rightarrow n}
    :=
    \left(L_{on}^{(t)}\right)^\top
    B_o
    L_{on}^{(t)}.
\end{equation}
Then
\begin{equation}
    R_{uu,o\rightarrow n}
    =
    B_{o\rightarrow n}\otimes G.
    \label{eq:app_Ruu_transfer_kron}
\end{equation}
The new block likelihood contribution is
\begin{equation}
    \frac{1}{\sigma^2}A_n^\top A_n
    =
    \frac{1}{\sigma^2}
    (T_n^\top T_n)\otimes G.
\end{equation}
Therefore the updated temporal likelihood statistic is
\begin{equation}
    B_n
    =
    B_{o\rightarrow n}
    +
    \frac{1}{\sigma^2}T_n^\top T_n.
    \label{eq:app_Bn_after_transfer}
\end{equation}
The resulting inducing precision has the Sylvester-compatible form
\begin{equation}
    \Lambda_{uu,n}
    =
    \left(K_{nn}^{(t)}\right)^{-1}\otimes K_s^{-1}
    +
    B_n\otimes G.
    \label{eq:app_Lambda_un_sylvester_form}
\end{equation}

Consequently, for \(q_*=\vecop(Q)\) and
\(z=\Lambda_{uu,n}^{-1}q_*=\vecop(Z)\), the system
\begin{equation}
    \Lambda_{uu,n}z=q_*
\end{equation}
is equivalent to the Sylvester-type equation
\begin{equation}
    K_s^{-1}Z
    \left(K_{nn}^{(t)}\right)^{-1}
    +
    GZB_n
    =
    Q,
    \label{eq:app_sylvester_after_transfer}
\end{equation}
assuming the vectorization convention used in the main text and symmetric
\(B_n\). Thus the SSGP-style old-likelihood-ratio transfer can preserve
the Sylvester solver under the temporal-only changing-basis and fixed
spatial-geometry assumptions.


\subsection{Explicit Sylvester Solver and Fast Quadratic Forms}
\label{app:explicit_sylvester_solver}

This subsection spells out the direct solver used whenever a vector of the form
\(D_u^{-1}q\) or a quadratic form \(q^\top D_u^{-1}q\) is required.  The derivation
is useful for three operations in the main method: posterior recovery,
predictive uncertainty, and the construction of the Schur complement.

Under the single-Kronecker invariant established above, the lower-right
inducing precision block has the form
\begin{equation}
    D_u
    =
    K_t^{-1}\otimes K_s^{-1}+B_n\otimes G,
    \label{eq:app_explicit_du}
\end{equation}
where \(K_t=K_{nn}^{(t)}\), \(K_s=K_{ZZ}^{(s)}\), and \(G=C^\top C\).
For a right-hand side \(q\in\mathbb{R}^{M_tM_s}\), write
\begin{equation}
    q=\vecop(Q),
    \qquad
    z=D_u^{-1}q=\vecop(Z),
    \qquad
    Q,Z\in\mathbb{R}^{M_s\times M_t}.
    \label{eq:app_qz_matrix_form}
\end{equation}
Using \((A\otimes B)\vecop(X)=\vecop(BXA^\top)\), the linear system
\(D_uz=q\) is equivalent to
\begin{equation}
    K_s^{-1}ZK_t^{-T}+GZB_n^\top=Q.
    \label{eq:app_general_sylvester_solver}
\end{equation}
When \(K_t\), \(K_s\), \(B_n\), and \(G\) are symmetric, this reduces to the
Sylvester-type equation
\begin{equation}
    K_s^{-1}ZK_t^{-1}+GZB_n=Q.
    \label{eq:app_symmetric_sylvester_solver}
\end{equation}
Thus the full \(M_tM_s\)-dimensional solve is replaced by a matrix equation over
an \(M_s\times M_t\) unknown.

\paragraph{Whitening.}
Introduce the whitened variables
\begin{equation}
    \widetilde Z=K_s^{-1/2}ZK_t^{-1/2},
    \qquad
    \widetilde Q=K_s^{1/2}QK_t^{1/2},
    \label{eq:app_whitened_ZQ}
\end{equation}
and the whitened likelihood factors
\begin{equation}
    \widetilde G=K_s^{1/2}GK_s^{1/2},
    \qquad
    \widetilde B=K_t^{1/2}B_nK_t^{1/2}.
    \label{eq:app_whitened_GB}
\end{equation}
Multiplying \eqref{eq:app_symmetric_sylvester_solver} on the left by
\(K_s^{1/2}\) and on the right by \(K_t^{1/2}\) gives
\begin{equation}
    \widetilde Z+\widetilde G\widetilde Z\widetilde B=\widetilde Q.
    \label{eq:app_whitened_sylvester_solver}
\end{equation}

\paragraph{Diagonal solution.}
Let
\begin{equation}
    \widetilde G=U_s\Gamma U_s^\top,
    \qquad
    \widetilde B=U_tH U_t^\top,
    \label{eq:app_eigendecomp_GB}
\end{equation}
where
\begin{equation}
    \Gamma=\diag(\gamma_1,\ldots,\gamma_{M_s}),
    \qquad
    H=\diag(\eta_1,\ldots,\eta_{M_t}).
\end{equation}
Define the rotated variables
\begin{equation}
    \widehat Z=U_s^\top\widetilde ZU_t,
    \qquad
    \widehat Q=U_s^\top\widetilde Q U_t.
    \label{eq:app_rotated_ZQ}
\end{equation}
Substitution into \eqref{eq:app_whitened_sylvester_solver} yields
\begin{equation}
    \widehat Z+\Gamma\widehat ZH=\widehat Q.
    \label{eq:app_rotated_sylvester}
\end{equation}
Because \(\Gamma\) and \(H\) are diagonal, the equation decouples entrywise:
\begin{equation}
    (1+\gamma_i\eta_j)\widehat Z_{ij}=\widehat Q_{ij},
    \qquad
    \widehat Z_{ij}=\frac{\widehat Q_{ij}}{1+\gamma_i\eta_j}.
    \label{eq:app_entrywise_solve}
\end{equation}
The full solution of the original system is recovered by reversing the transforms:
\begin{equation}
    \widetilde Z=U_s\widehat ZU_t^\top,
    \qquad
    Z=K_s^{1/2}\widetilde ZK_t^{1/2},
    \qquad
    z=\vecop(Z).
    \label{eq:app_inverse_transforms}
\end{equation}
The same eigensystems can be reused for all right-hand sides in a block.  In
particular, the solves needed for \(D_u^{-1}h_u\), \(D_u^{-1}B_{\beta u}^\top\),
and \(D_u^{-1}a_*\) differ only in the matrix \(Q\).

\paragraph{Quadratic form without reconstructing the solution.}
For prediction, one often needs only the scalar posterior-uncertainty correction
\(q^\top D_u^{-1}q\).  Since \(q=\vecop(Q)\) and \(z=\vecop(Z)\),
\begin{equation}
    q^\top D_u^{-1}q=q^\top z=\langle Q,Z\rangle_F.
    \label{eq:app_qform_matrix_inner}
\end{equation}
This Frobenius inner product is invariant under the whitening and orthogonal
rotations above:
\begin{equation}
    \langle Q,Z\rangle_F
    =\langle \widetilde Q,\widetilde Z\rangle_F
    =\langle \widehat Q,\widehat Z\rangle_F.
    \label{eq:app_inner_invariance}
\end{equation}
Indeed, using \(Q=K_s^{-1/2}\widetilde QK_t^{-1/2}\) and
\(Z=K_s^{1/2}\widetilde ZK_t^{1/2}\), cyclic invariance of the trace gives
\begin{equation}
    \tr(Q^\top Z)
    =\tr(\widetilde Q^\top \widetilde Z),
\end{equation}
and orthogonality of \(U_s,U_t\) gives
\begin{equation}
    \tr(\widetilde Q^\top\widetilde Z)
    =\tr(\widehat Q^\top\widehat Z).
\end{equation}
Combining this identity with \eqref{eq:app_entrywise_solve} yields the exact
fast quadratic form
\begin{equation}
    q^\top D_u^{-1}q
    =
    \sum_{i=1}^{M_s}\sum_{j=1}^{M_t}
    \frac{\widehat Q_{ij}^{2}}{1+\gamma_i\eta_j}.
    \label{eq:app_fast_quadratic_form}
\end{equation}
Therefore, when only the scalar predictive uncertainty contribution is needed,
one does not need to reconstruct \(\widehat Z\to\widetilde Z\to Z\to z\).  The
scalar can be accumulated directly in the diagonalized coordinates.  When the
full vector \(D_u^{-1}q\) is needed, as in Schur-complement recovery, the inverse
transforms in \eqref{eq:app_inverse_transforms} are applied.

\paragraph{Complexity.}
A dense direct solve with the mixed inducing precision costs
\(\mathcal{O}((M_tM_s)^3)\).  The diagonalized Sylvester solver costs
\begin{equation}
    \mathcal{O}\!\big(M_t^3+M_s^3+M_t^2M_s+M_tM_s^2\big),
    \label{eq:app_sylvester_complexity}
\end{equation}
up to constant factors.  The two cubic terms come from eigendecompositions of
the whitened temporal and spatial factors, and the mixed terms come from the
left and right basis transforms.  Since the spatial inducing geometry is fixed,
the spatial whitening and eigendecomposition can be cached across online blocks,
so the amortized online cost is dominated by temporal updates and two-sided
matrix transforms.

\subsection{\texorpdfstring{When \(R_{uu,o}=B_o\otimes G\) Can Fail}{When Ruu,o = Bo tensor G Can Fail}}
\label{app:failure_cases_Ro_kron}

The single-Kronecker representation \(R_{uu,o}=B_o\otimes G\) is not
unconditional. It may fail in the following cases.

\paragraph{Changing spatial observation pattern.}
If the spatial observation pattern or spatial mini-batch changes across
blocks, then the spatial projection matrix may become block-dependent:
\begin{equation}
    C_j\neq C.
\end{equation}
Consequently,
\begin{equation}
    G_j=C_j^\top C_j
\end{equation}
also changes. The old $u$--$u$ likelihood precision becomes
\begin{equation}
    R_{uu,o}
    =
    \sum_{j<n}
    B_j\otimes G_j,
\end{equation}
rather than a single Kronecker product \(B_o\otimes G\). This is still a
sum-of-Kronecker structure, but the simple closed-form Sylvester equation
in \eqref{eq:app_sylvester_after_transfer} no longer applies directly.
One would need a generalized Sylvester solver or an iterative method using
Kronecker matrix-vector products.

\paragraph{Changing spatial inducing locations.}
If the spatial inducing locations \(Z_s\) are updated online, then
\begin{equation}
    K_s=K_{ZZ}^{(s)},
    \qquad
    C=K_{XZ}^{(s)}\left(K_{ZZ}^{(s)}\right)^{-1},
    \qquad
    G=C^\top C
\end{equation}
all change over time. In that case, the old-to-new transfer operator need
not have the form \(L_{on}^{(t)}\otimes I_s\), and the likelihood
precision is generally not a single \(B_o\otimes G\) term.

\paragraph{Recovering \(R_{uu,o}\) from a dense posterior covariance.}
If one stores a full joint posterior over \((\beta,u_o)\) and then uses only
the marginal covariance of \(u_o\), the inverse of that marginal covariance
does not generally equal the lower-right precision block of the joint
posterior. The structure-preserving statistic should instead be interpreted in
information form:
\begin{equation}
    R_{uu,o}=D_{u,o}-K_{oo}^{-1},
\end{equation}
where \(D_{u,o}\) is the lower-right posterior precision block. This is a
single Kronecker product only if that precision block preserves the required
structure. For the Sylvester-preserving implementation, the old $u$--$u$
likelihood natural precision should be maintained directly as \(B_o\otimes G\),
rather than recovered from an arbitrary dense covariance.

\paragraph{Non-Gaussian likelihoods or non-conjugate approximations.}
For Gaussian likelihoods, the likelihood precision contribution is
\begin{equation}
    \frac{1}{\sigma^2}A_j^\top A_j.
\end{equation}
With non-Gaussian likelihoods, approximate inference methods such as
Laplace, expectation propagation, or stochastic variational inference may
produce likelihood-site precisions that are not simply proportional to
\(A_j^\top A_j\), or that vary by observation. The resulting precision
may fail to share a common spatial factor \(G\).


\paragraph{Possible extension: sum-of-Kronecker Route B.}
A natural extension, not used in the main method, is to maintain the old
$u$--$u$ likelihood precision as a sum of Kronecker products:
\begin{equation}
    R_{uu,o}
    =
    \sum_{r=1}^{R_o} B_r\otimes G_r .
    \label{eq:app_sum_kron_ruu_old}
\end{equation}
This representation is more general than the single-Kronecker invariant
\(B_o\otimes G\).  For example, it can represent changing spatial
observation masks by assigning a separate spatial factor
\(G_r=C_r^\top C_r\) to each observation pattern, or to a compressed group
of blocks sharing the same pattern.

Under the temporal-only changing-basis map used in the main text,
\(L_{on}=L_t\otimes I_s\), the transferred old likelihood precision becomes
\begin{equation}
    R_{uu,o\to n}
    =
    \sum_{r=1}^{R_o}
    \left(L_t^\top B_r L_t\right)\otimes G_r .
    \label{eq:app_sum_kron_transfer}
\end{equation}
The new block adds one further Kronecker term,
\begin{equation}
    R_{uu,n}^{\mathrm{new}}
    =
    B_n\otimes G_n,
    \label{eq:app_sum_kron_new_block}
\end{equation}
where \(G_n=C_n^\top C_n\) may differ from previous spatial factors if the
current spatial observation pattern is different.  If the spatial inducing
basis also changes and the old-to-new map remains separable,
\(L_{on}=L_t\otimes L_s\), then the same idea gives
\begin{equation}
    R_{uu,o\to n}
    =
    \sum_{r=1}^{R_o}
    \left(L_t^\top B_r L_t\right)
    \otimes
    \left(L_s^\top G_r L_s\right).
    \label{eq:app_sum_kron_spatial_transfer}
\end{equation}
This covers only structured spatial-basis changes.  If the old-to-new map is
not separable, then the sum-of-Kronecker form need not be preserved.

With the sum-of-Kronecker representation, the $u$-block precision has the
form
\begin{equation}
    D_u
    =
    K_t^{-1}\otimes K_s^{-1}
    +
    \sum_{r=1}^{R} B_r\otimes G_r,
    \label{eq:app_sum_kron_du}
\end{equation}
where the sum includes the transferred old terms and the current block term.
The corresponding matrix equation for \(v=\vecop(M)\) is no longer the
single Sylvester equation \eqref{eq:app_sylvester_after_transfer}; instead it
becomes
\begin{equation}
    K_s^{-1}M K_t^{-1}
    +
    \sum_{r=1}^{R} G_r M B_r
    =
    H,
    \label{eq:app_sum_kron_matrix_equation}
\end{equation}
assuming the same vectorization convention as above and symmetric temporal
factors.  This system can be solved by an iterative method such as conjugate
gradients, using Kronecker matrix--vector products rather than materializing
\(D_u\).  In particular,
\begin{equation}
    \left(B_r\otimes G_r\right)\vecop(M)
    =
    \vecop\!\left(G_r M B_r^\top\right).
    \label{eq:app_sum_kron_mvm}
\end{equation}
The Route~B Schur-complement recovery can then still be used in principle:
every occurrence of \(D_u^{-1}\) in the posterior mean, covariance, and
predictive-variance formulas would be evaluated by the iterative
sum-of-Kronecker solver.  However, this extension loses the simple direct
Sylvester solve, increases storage because multiple \((B_r,G_r)\) factors
must be retained or compressed, and requires preconditioning for large
streams.  We therefore treat \eqref{eq:app_sum_kron_ruu_old}--\eqref{eq:app_sum_kron_mvm}
as a future extension rather than as part of the main Route~B method.

\paragraph{Practical implication.}
To retain the simple Sylvester solver, the scalable version of the method
should explicitly assume fixed spatial inducing geometry and should
maintain the old $u$--$u$ likelihood precision as a Kronecker natural statistic:
\begin{equation}
    R_{uu,o}=B_o\otimes G.
\end{equation}
When these assumptions are violated, the model can still be implemented,
but the solver must be replaced by a more general structured linear
algebra routine, such as a sum-of-Kronecker conjugate-gradient method or a
generalized Sylvester solver.

\clearpage
\section{Formula-by-Formula English Explanation and Speaking Notes}
This appendix is an annotated companion to the theory document. It keeps the main formulas in their original order and adds a plain-English derivation guide for discussion with a supervisor. The original source TeX file is not modified; this is a generated annotated copy.

\subsection{1. Observation Model}
\paragraph{Formula 1.1 Scalar observation}
\[
y(t_n,s_i)=\phi(t_n,s_i)^\top\beta+f(t_n,s_i)+\epsilon_{n,i},
\qquad
\epsilon_{n,i}\sim\mathcal N(0,\sigma^2).
\]
\textbf{English explanation.}
\begin{enumerate}[leftmargin=2em]
\item \(y(t_n,s_i)\) is the scalar observation at time \(t_n\) and spatial location \(s_i\).
\item \(\phi(t_n,s_i)^\top\beta\) is the parametric linear trend captured by engineered features.
\item \(f(t_n,s_i)\) is the nonparametric GP residual that captures smooth spatio-temporal structure not explained by the linear trend.
\item \(\epsilon_{n,i}\) is Gaussian observation noise with variance \(\sigma^2\).
\item The key modeling choice is joint observation: the linear trend and GP residual explain \(y\) together.
\end{enumerate}
\textbf{Speaking point.} Route B starts from a joint likelihood for the raw observation \(y\), not from stale residuals produced by an old \(\hat\beta\).

\paragraph{Formula 1.2 Block observation}
\[
\mathbf y_n=\Phi_n\beta+\mathbf f_n+\epsilon_n,
\qquad
\epsilon_n\sim\mathcal N(0,\sigma^2 I_{B_n}).
\]
\textbf{English explanation.}
\begin{enumerate}[leftmargin=2em]
\item All observations in online block \(n\) are stacked into \(\mathbf y_n\).
\item \(\Phi_n\) stacks the linear features row by row.
\item \(\mathbf f_n\) stacks the latent GP residual values at the same points.
\item The noise vector is independent Gaussian noise with covariance \(\sigma^2 I\).
\item This matrix form lets us write a Gaussian likelihood and derive natural-parameter updates.
\end{enumerate}

\subsection{2. Separable GP and Kronecker Prior}
\paragraph{Formula 2.1 Separable spatio-temporal kernel}
\[
f\sim\mathcal{GP}(0,k),
\qquad
k((s,t),(s',t'))=k_s(s,s')k_t(t,t').
\]
\textbf{English explanation.}
\begin{enumerate}[leftmargin=2em]
\item \(f\) is a zero-mean Gaussian process.
\item The kernel \(k\) measures covariance between two spatio-temporal inputs.
\item The separable assumption factorizes covariance into a spatial part and a temporal part.
\item This is the structural assumption that enables Kronecker algebra.
\end{enumerate}

\paragraph{Formula 2.2 Full-grid covariance}
\[
K_{ff}=K_t\otimes K_s.
\]
\textbf{English explanation.}
\begin{enumerate}[leftmargin=2em]
\item On a full time-space grid, \(K_t\) stores temporal covariance and \(K_s\) stores spatial covariance.
\item The separable kernel makes the full covariance a Kronecker product.
\item The Kronecker product represents a large covariance using two smaller factors.
\item This structure is what makes the later Sylvester solver possible.
\end{enumerate}

\subsection{3. Temporal HiPPO-RFF Inducing Representation}
\paragraph{Formula 3.1 Scaled Legendre basis}
\[
g_\ell^{(T)}(x)=\sqrt{\frac{2\ell+1}{T}}\,P_\ell\!\left(\frac{2x}{T}-1\right),
\qquad x\in[0,T].
\]
\textbf{English explanation.}
\begin{enumerate}[leftmargin=2em]
\item \(P_\ell\) is the \(\ell\)-th Legendre polynomial on \([-1,1]\).
\item The transformation \(2x/T-1\) maps the current time interval \([0,T]\) to \([-1,1]\).
\item The square-root factor normalizes the basis on the interval.
\item Since \(T\) changes online, the inducing coordinates also change.
\end{enumerate}

\paragraph{Formula 3.2 Temporal interdomain inducing variable}
\[
u_\ell^{(t)}(T)=\int_0^T g_\ell^{(T)}(x)f_t(x)\,dx.
\]
\textbf{English explanation.}
\begin{enumerate}[leftmargin=2em]
\item This inducing variable is not a point evaluation. It is a projection of the temporal GP onto basis \(g_\ell^{(T)}\).
\item The integral measures how much of the temporal function lies in that basis direction.
\item Therefore it is an interdomain inducing variable.
\item Changing \(T\) changes the coordinate system.
\end{enumerate}

\paragraph{Formula 3.3 RFF approximation and basis matrix}
\[
k_t(\tau,\tau')\approx
\sum_r w_r[
\cos(\omega_r\tau)\cos(\omega_r\tau')
+
\sin(\omega_r\tau)\sin(\omega_r\tau')
].
\]
\[
\big[B_t(T)\big]_{\ell,r}=\sqrt{w_r}I_\ell^{\cos}(\omega_r;T),
\qquad
\big[B_t(T)\big]_{\ell,R+r}=\sqrt{w_r}I_\ell^{\sin}(\omega_r;T).
\]
\textbf{English explanation.}
\begin{enumerate}[leftmargin=2em]
\item The temporal kernel is approximated by random Fourier features.
\item Each frequency contributes a cosine and sine component.
\item \(B_t(T)\) maps the HiPPO-Legendre inducing basis into the RFF feature space.
\item The entries are oscillatory integrals between the Legendre basis and Fourier features.
\item This matrix lets us compute temporal inducing covariances efficiently.
\end{enumerate}

\paragraph{Formula 3.4 Temporal inducing covariance and cross-covariance}
\[
K_{uu}^{(t)}(T)\approx B_t(T)B_t(T)^\top,
\qquad
k_{fu}^{(t)}(\tau;T)\approx \varphi_t(\tau)^\top B_t(T)^\top.
\]
\textbf{English explanation.}
\begin{enumerate}[leftmargin=2em]
\item \(K_{uu}^{(t)}\) is the prior covariance among temporal inducing variables.
\item In the RFF representation, it is approximated by \(B_tB_t^\top\).
\item \(k_{fu}^{(t)}\) is the covariance between a temporal input \(\tau\) and the inducing variables.
\item It is computed by projecting the RFF feature vector onto the inducing basis.
\end{enumerate}

\subsection{4. Mixed Spatio-Temporal Inducing Variables}
\paragraph{Formula 4.1 Mixed inducing variable}
\[
u_{\ell,j}(T)=\int_0^T g_\ell^{(T)}(x)f(x,z_j^s)\,dx.
\]
\textbf{English explanation.}
\begin{enumerate}[leftmargin=2em]
\item \(\ell\) indexes the temporal basis.
\item \(j\) indexes the spatial inducing location.
\item The GP is evaluated at a fixed spatial inducing point and integrated over time.
\item The resulting variable combines temporal interdomain structure and spatial inducing points.
\end{enumerate}

\paragraph{Formula 4.2 Kronecker inducing covariance}
\[
K_{uu}^{st}=K_{uu}^{(t)}\otimes K_{ZZ}^{(s)},
\qquad
K_{fu,n}^{st}=K_{fu,n}^{(t)}\otimes K_{XZ}^{(s)}.
\]
\textbf{English explanation.}
\begin{enumerate}[leftmargin=2em]
\item \(K_{uu}^{st}\) is the prior covariance of all mixed inducing variables.
\item The temporal factor is \(K_{uu}^{(t)}\), and the spatial factor is \(K_{ZZ}^{(s)}\).
\item Separability makes the covariance a Kronecker product.
\item The observation-inducing cross-covariance factorizes in the same way.
\end{enumerate}

\paragraph{Formula 4.3 Sparse projection matrix}
\[
A_n=K_{fu,n}^{st}(K_{uu,n}^{st})^{-1}
=
\left[K_{fu,n}^{(t)}(K_{uu,n}^{(t)})^{-1}\right]
\otimes
\left[K_{XZ}^{(s)}(K_{ZZ}^{(s)})^{-1}\right].
\]
\[
T_n:=K_{fu,n}^{(t)}(K_{uu,n}^{(t)})^{-1},
\qquad
C:=K_{XZ}^{(s)}(K_{ZZ}^{(s)})^{-1},
\qquad
A_n=T_n\otimes C.
\]
\textbf{English explanation.}
\begin{enumerate}[leftmargin=2em]
\item \(A_n\) is the sparse GP conditional-mean projection from inducing variables to observations.
\item It has the standard sparse GP form \(K_{fu}K_{uu}^{-1}\).
\item Since both covariance matrices are Kronecker-structured, the projection factorizes.
\item \(T_n\) handles temporal projection; \(C\) handles spatial projection.
\item The identity \(A_n=T_n\otimes C\) is the computational backbone of Route B.
\end{enumerate}

\subsection{5. Why Residual-Based Online Learning Fails}
\paragraph{Formula 5.1 Stale residual identity}
\[
r_j(\hat\beta_n)
=y_j-\Phi_j\hat\beta_n
=y_j-\Phi_j\hat\beta_{n-1}-\Phi_j(\hat\beta_n-\hat\beta_{n-1})
=r_j(\hat\beta_{n-1})-\Phi_j(\hat\beta_n-\hat\beta_{n-1}).
\]
\textbf{English explanation.}
\begin{enumerate}[leftmargin=2em]
\item \(r_j(\hat\beta_n)\) is the residual that historical block \(j\) would have under the latest linear coefficient.
\item The first equality is the residual definition.
\item The second equality decomposes the latest coefficient into the old coefficient plus a change.
\item The last equality shows that the old residual must be corrected when \(\beta\) changes.
\item If old raw observations are unavailable, residual-based online learning cannot repair the old targets.
\end{enumerate}

\paragraph{Formula 5.2 Joint likelihood avoids stale residuals}
\[
y_j\mid\beta,u_j\sim
\mathcal N(\Phi_j\beta+A_j u_j,\sigma^2 I).
\]
\textbf{English explanation.}
\begin{enumerate}[leftmargin=2em]
\item Route B does not train the GP on residuals computed from a point estimate of \(\beta\).
\item It keeps a joint likelihood for the raw observation \(y_j\).
\item Historical data constrain \(\beta\) and \(u\) jointly, avoiding stale residual targets.
\end{enumerate}

\subsection{6. Joint ELBO and the Missing Mean-Field Cross Term}
\paragraph{Formula 6.1 Structured joint posterior}
\[
q_n(\beta,u_n)
=
\mathcal N
\left(
\begin{bmatrix}\beta\\u_n\end{bmatrix}
\middle|
\begin{bmatrix}m_{\beta,n}\\m_{u,n}\end{bmatrix},
S_n
\right),
\qquad
S_n^{-1}=\Lambda_n.
\]
\textbf{English explanation.}
\begin{enumerate}[leftmargin=2em]
\item The approximate posterior is a joint Gaussian over \((\beta,u_n)\), not a product of marginals.
\item The mean contains both the linear coefficient mean and inducing mean.
\item The covariance allows beta-u posterior dependence.
\item The precision form \(\Lambda_n=S_n^{-1}\) is convenient for Gaussian updates.
\end{enumerate}

\paragraph{Formula 6.2 Expected log likelihood}
\[
\begin{aligned}
\mathbb E_q[\log p(y_n\mid\beta,u_n)]
=&-\frac{B_n}{2}\log(2\pi\sigma^2)\\
&-\frac{1}{2\sigma^2}
\Big[
\|y_n-\Phi_n m_\beta-A_n m_u\|^2
+\operatorname{tr}(\Phi_nS_{\beta\beta}\Phi_n^\top)\\
&\qquad
+\operatorname{tr}(A_nS_{uu}A_n^\top)
+2\operatorname{tr}(\Phi_nS_{\beta u}A_n^\top)
\Big].
\end{aligned}
\]
\textbf{English explanation.}
\begin{enumerate}[leftmargin=2em]
\item The first term is the Gaussian likelihood normalization.
\item The squared residual uses the posterior means of \(\beta\) and \(u\).
\item The \(\beta\)-trace term is output uncertainty due to the linear coefficient.
\item The \(u\)-trace term is output uncertainty due to the GP inducing variables.
\item The cross trace is the covariance contribution between the linear trend and GP residual.
\item Mean-field removes this term by setting \(S_{\beta u}=0\).
\end{enumerate}
\textbf{Speaking point.} The direct theoretical difference between Route B and mean-field is the retained beta-u covariance term.

\subsection{7. Fixed-Basis Exact Gaussian Update}
\paragraph{Formula 7.1 Joint state and likelihood matrix}
\[
z=\begin{bmatrix}\beta\\u\end{bmatrix},
\qquad
H_n=\begin{bmatrix}\Phi_n&A_n\end{bmatrix},
\qquad
y_n\mid z\sim\mathcal N(H_nz,\sigma^2I).
\]
\textbf{English explanation.}
\begin{enumerate}[leftmargin=2em]
\item \(z\) concatenates the linear coefficients and inducing variables.
\item \(H_n\) concatenates the linear and GP design matrices.
\item \(H_nz\) reproduces the joint observation mean.
\item With a fixed basis, the same state can be updated by standard Gaussian conditioning.
\end{enumerate}

\paragraph{Formula 7.2 Natural-parameter update}
\[
\Lambda_n=\Lambda_{n-1}+\sigma^{-2}H_n^\top H_n,
\qquad
h_n=h_{n-1}+\sigma^{-2}H_n^\top y_n.
\]
\[
S_n=\Lambda_n^{-1},
\qquad
m_n=S_nh_n.
\]
\textbf{English explanation.}
\begin{enumerate}[leftmargin=2em]
\item A Gaussian likelihood contributes \(\sigma^{-2}H^\top H\) to precision.
\item It contributes \(\sigma^{-2}H^\top y\) to the information vector.
\item The new natural parameters are old natural parameters plus current likelihood contributions.
\item The covariance is the inverse precision.
\item The mean is recovered by \(m=\Lambda^{-1}h\).
\end{enumerate}

\paragraph{Formula 7.3 Block expansion}
\[
H_n^\top H_n
=
\begin{bmatrix}
\Phi_n^\top\Phi_n & \Phi_n^\top A_n\\
A_n^\top\Phi_n & A_n^\top A_n
\end{bmatrix}.
\]
\textbf{English explanation.}
\begin{enumerate}[leftmargin=2em]
\item The top-left block is information about \(\beta\).
\item The bottom-right block is information about \(u\).
\item The off-diagonal blocks encode coupling between the linear trend and GP residual.
\item Route B keeps these cross blocks.
\end{enumerate}

\subsection{8. Changing-Basis Transfer}
\paragraph{Formula 8.1 Old and new inducing variables are different}
\[
u_\ell^{old}=\int_0^{T_{n-1}}g_\ell^{(T_{n-1})}(x)f(x)\,dx,
\qquad
u_\ell^{new}=\int_0^{T_n}g_\ell^{(T_n)}(x)f(x)\,dx.
\]
\textbf{English explanation.}
\begin{enumerate}[leftmargin=2em]
\item The old inducing variable is defined using the old temporal basis and horizon.
\item The new inducing variable is defined using the new basis and horizon.
\item They are different random variables even if they share the same index.
\item Therefore old posterior information must be transferred to the new coordinates.
\end{enumerate}

\paragraph{Formula 8.2 Streaming old-likelihood ratio}
\[
\frac{q_{n-1}(\beta,u_o)}{p(\beta)p(u_o)}.
\]
\textbf{English explanation.}
\begin{enumerate}[leftmargin=2em]
\item \(q_{n-1}\) is the posterior after old data.
\item \(p(\beta)p(u_o)\) is the base prior in the old coordinates.
\item Their ratio isolates the information contributed by old observations.
\item Streaming sparse GP transfers this likelihood ratio, not the posterior itself.
\end{enumerate}

\paragraph{Formula 8.3 Route B streaming VFE update}
\[
q_n(\beta,u_n)\propto
p(\beta)p(u_n)
\exp\left\{
\mathbb E_{p(u_o\mid u_n)}
\left[
\log\frac{q_{n-1}(\beta,u_o)}{p(\beta)p(u_o)}
\right]
+\log p(y_n\mid\beta,u_n)
\right\}.
\]
\textbf{English explanation.}
\begin{enumerate}[leftmargin=2em]
\item The left side is the new posterior after block \(n\).
\item \(p(\beta)p(u_n)\) is the base prior in the new coordinates.
\item The first exponential term is historical likelihood information written in old coordinates.
\item The expectation under \(p(u_o\mid u_n)\) transfers that information to the new coordinates.
\item The second term is the current block likelihood.
\item This is the core changing-basis Route B update.
\end{enumerate}

\subsection{9. Old Likelihood Natural Parameters}
\paragraph{Formula 9.1 Natural-parameter representation}
\[
\log\frac{q_{n-1}(\beta,u_o)}{p(\beta)p(u_o)}
=
-\frac12
\begin{bmatrix}\beta\\u_o\end{bmatrix}^\top
R_o^z
\begin{bmatrix}\beta\\u_o\end{bmatrix}
+
\begin{bmatrix}\beta\\u_o\end{bmatrix}^\top r_o^z
+c_o.
\]
\textbf{English explanation.}
\begin{enumerate}[leftmargin=2em]
\item The old likelihood ratio is a Gaussian factor, so its log is quadratic.
\item \(R_o^z\) is the likelihood precision contributed by old data.
\item \(r_o^z\) is the corresponding information vector.
\item Constants can be ignored for posterior recovery.
\end{enumerate}

\paragraph{Formula 9.2 Block structure}
\[
R_o^z=
\begin{bmatrix}
R_{\beta\beta,o}&R_{\beta u,o}\\
R_{u\beta,o}&R_{uu,o}
\end{bmatrix},
\qquad
r_o^z=
\begin{bmatrix}
r_{\beta,o}\\r_{u,o}
\end{bmatrix}.
\]
\textbf{English explanation.}
\begin{enumerate}[leftmargin=2em]
\item \(R_{\beta\beta,o}\) is old information about \(\beta\).
\item \(R_{uu,o}\) is old information about \(u\).
\item \(R_{\beta u,o}\) stores beta-u coupling information.
\item Route B explicitly retains this block.
\end{enumerate}

\subsection{10. Closed-Form Transfer of Old Information}
\paragraph{Formula 10.1 GP conditional from new to old inducing variables}
\[
p(u_o\mid u_n)=\mathcal N(L_{on}u_n,\Sigma_{o\mid n}),
\qquad
L_{on}=K_{on}K_{nn}^{-1}.
\]
\textbf{English explanation.}
\begin{enumerate}[leftmargin=2em]
\item The GP prior defines a joint Gaussian between old and new inducing variables.
\item The conditional mean predicts old coordinates from new coordinates.
\item \(L_{on}=K_{on}K_{nn}^{-1}\) is the standard Gaussian conditioning operator.
\item The conditional covariance affects only constants in the likelihood transfer.
\end{enumerate}

\paragraph{Formula 10.2 Deterministic transfer matrix}
\[
M_{on}:=
\begin{bmatrix}
I_d&0\\
0&L_{on}
\end{bmatrix}.
\]
\textbf{English explanation.}
\begin{enumerate}[leftmargin=2em]
\item \(\beta\) is unchanged across bases, so it uses the identity map.
\item \(u_o\) is predicted from \(u_n\), so it uses \(L_{on}\).
\item The block matrix maps new coordinates into the old likelihood coordinate system.
\end{enumerate}

\paragraph{Formula 10.3 Compact transfer}
\[
R_{o\rightarrow n}^z=M_{on}^\top R_o^zM_{on},
\qquad
r_{o\rightarrow n}^z=M_{on}^\top r_o^z.
\]
\textbf{English explanation.}
\begin{enumerate}[leftmargin=2em]
\item A quadratic form transforms as \(R\mapsto M^\top RM\).
\item A linear information vector transforms as \(r\mapsto M^\top r\).
\item This is simply Gaussian natural-parameter change of coordinates.
\end{enumerate}

\paragraph{Formula 10.4 Block transfer equations}
\[
\begin{aligned}
R_{\beta\beta,o\rightarrow n}&=R_{\beta\beta,o},\\
R_{\beta u,o\rightarrow n}&=R_{\beta u,o}L_{on},\\
R_{u\beta,o\rightarrow n}&=L_{on}^\top R_{u\beta,o},\\
R_{uu,o\rightarrow n}&=L_{on}^\top R_{uu,o}L_{on},\\
r_{\beta,o\rightarrow n}&=r_{\beta,o},\\
r_{u,o\rightarrow n}&=L_{on}^\top r_{u,o}.
\end{aligned}
\]
\textbf{English explanation.}
\begin{enumerate}[leftmargin=2em]
\item The beta-beta block is unchanged because beta is unchanged.
\item The beta-u block transforms on the \(u\) side by right multiplication with \(L_{on}\).
\item The u-beta block transforms symmetrically by left multiplication with \(L_{on}^\top\).
\item The u-u block transforms on both sides.
\item The beta information vector is unchanged.
\item The u information vector is pulled into the new coordinates by \(L_{on}^\top\).
\end{enumerate}

\subsection{11. New Block Assimilation}
\paragraph{Formula 11.1 New likelihood natural parameters}
\[
H_n=\begin{bmatrix}\Phi_n&A_n\end{bmatrix},
\qquad
R_{\mathrm{new},n}^z=\frac1{\sigma^2}H_n^\top H_n,
\qquad
r_{\mathrm{new},n}^z=\frac1{\sigma^2}H_n^\top y_n.
\]
\textbf{English explanation.}
\begin{enumerate}[leftmargin=2em]
\item \(H_n\) combines the linear and GP design matrices.
\item A Gaussian likelihood contributes \(H^\top H/\sigma^2\) to precision.
\item It contributes \(H^\top y/\sigma^2\) to the information vector.
\item The target is raw \(y_n\), not a precomputed residual.
\end{enumerate}

\paragraph{Formula 11.2 Accumulated likelihood statistics}
\[
\begin{aligned}
R_{\beta\beta,n}
&=R_{\beta\beta,o}+\frac1{\sigma^2}\Phi_n^\top\Phi_n,\\
R_{\beta u,n}
&=R_{\beta u,o}L_{on}+\frac1{\sigma^2}\Phi_n^\top A_n,\\
R_{u\beta,n}
&=L_{on}^\top R_{u\beta,o}+\frac1{\sigma^2}A_n^\top\Phi_n,\\
R_{uu,n}
&=L_{on}^\top R_{uu,o}L_{on}+\frac1{\sigma^2}A_n^\top A_n,\\
r_{\beta,n}
&=r_{\beta,o}+\frac1{\sigma^2}\Phi_n^\top y_n,\\
r_{u,n}
&=L_{on}^\top r_{u,o}+\frac1{\sigma^2}A_n^\top y_n.
\end{aligned}
\]
\textbf{English explanation.}
\begin{enumerate}[leftmargin=2em]
\item Every line has the same pattern: transferred old information plus new block information.
\item The beta-beta statistic only accumulates \(\Phi^\top\Phi\).
\item The beta-u statistic transfers the old cross block and adds the new cross likelihood.
\item The u-beta block is the transpose-side counterpart.
\item The u-u statistic transfers the old GP precision and adds \(A^\top A\).
\item The beta information vector accumulates \(\Phi^\top y\).
\item The u information vector transfers old GP information and adds \(A^\top y\).
\end{enumerate}

\paragraph{Formula 11.3 Add priors to form posterior natural parameters}
\[
\begin{aligned}
\Lambda_{\beta\beta,n}&=P_\beta+R_{\beta\beta,n},\\
\Lambda_{\beta u,n}&=R_{\beta u,n},\\
\Lambda_{uu,n}&=K_{nn}^{-1}+R_{uu,n},\\
h_{\beta,n}&=P_\beta m_{\beta,0}+r_{\beta,n},\\
h_{u,n}&=r_{u,n}.
\end{aligned}
\]
\textbf{English explanation.}
\begin{enumerate}[leftmargin=2em]
\item Posterior precision equals prior precision plus likelihood precision.
\item The beta prior contributes \(P_\beta\).
\item The beta-u precision has no prior cross term because the base prior factorizes.
\item The inducing prior contributes \(K_{nn}^{-1}\).
\item The beta information vector receives \(P_\beta m_{\beta,0}\).
\item The inducing prior mean is zero, so \(h_u\) only contains likelihood information.
\end{enumerate}

\subsection{12. Projected-Prior Baseline}
\paragraph{Formula 12.1 Moment projection baseline}
\[
p_n^{\mathrm{proj}}(\beta,u_n)
=
\int p(u_n\mid u_o)q_{n-1}(\beta,u_o)\,du_o.
\]
\[
\begin{aligned}
m_{\beta,n|n-1}^{proj}&=m_{\beta,o},\\
m_{u,n|n-1}^{proj}&=L_{no}m_{u,o},\\
S_{\beta u,n|n-1}^{proj}&=S_{\beta u,o}L_{no}^\top,\\
S_{uu,n|n-1}^{proj}
&=K_{nn}+L_{no}(S_{uu,o}-K_{oo})L_{no}^\top.
\end{aligned}
\]
\textbf{English explanation.}
\begin{enumerate}[leftmargin=2em]
\item Projected-prior transfer propagates posterior moments into the next prior.
\item Beta remains unchanged.
\item The inducing mean is projected using a conditional mean operator.
\item The cross covariance can also be projected.
\item The u-u covariance follows GP conditional moment propagation.
\item This is a diagnostic ablation, not the main Route B old-likelihood-ratio method.
\end{enumerate}

\subsection{13. Schur Complement Posterior Recovery}
\paragraph{Formula 13.1 Precision block notation}
\[
\Lambda_n=
\begin{bmatrix}
A_\beta&B_{\beta u}\\
B_{\beta u}^\top&D_u
\end{bmatrix},
\qquad
h_n=
\begin{bmatrix}
h_\beta\\h_u
\end{bmatrix}.
\]
\textbf{English explanation.}
\begin{enumerate}[leftmargin=2em]
\item \(A_\beta\) is the small beta precision block.
\item \(D_u\) is the large inducing precision block.
\item \(B_{\beta u}\) stores beta-u coupling in precision form.
\item We avoid inverting the full matrix by using block algebra.
\end{enumerate}

\paragraph{Formula 13.2 Schur complement precision}
\[
\Lambda_{\beta|u}
:=
A_\beta-B_{\beta u}D_u^{-1}B_{\beta u}^\top.
\]
\textbf{English explanation.}
\begin{enumerate}[leftmargin=2em]
\item This is the effective beta precision after eliminating the inducing block.
\item The correction term accounts for coupling through the GP block.
\item It is a precision matrix, not a covariance.
\item Its inverse gives \(S_{\beta\beta}\).
\end{enumerate}

\paragraph{Formula 13.3 Required Sylvester solves}
\[
D_uW=B_{\beta u}^\top,
\qquad
D_uv_h=h_u,
\qquad
\Lambda_{\beta|u}=A_\beta-B_{\beta u}W.
\]
\textbf{English explanation.}
\begin{enumerate}[leftmargin=2em]
\item \(W\) is \(D_u^{-1}B_{\beta u}^\top\), computed by solving systems.
\item \(v_h\) is \(D_u^{-1}h_u\).
\item Since the beta dimension is small, only \(d\) right-hand sides are needed for \(W\).
\item Together with \(v_h\), posterior recovery requires \(d+1\) large-block solves.
\end{enumerate}

\paragraph{Formula 13.4 Posterior means}
\[
m_\beta=
\Lambda_{\beta|u}^{-1}(h_\beta-B_{\beta u}v_h),
\qquad
m_u=v_h-Wm_\beta.
\]
\textbf{English explanation.}
\begin{enumerate}[leftmargin=2em]
\item \(h_\beta-B_{\beta u}v_h\) is the effective beta information vector after eliminating \(u\).
\item Multiplying by the inverse Schur complement gives \(m_\beta\).
\item The expression for \(m_u\) comes from the second block row of the precision system.
\item It adjusts the inducing mean by the beta-u coupling.
\end{enumerate}

\paragraph{Formula 13.5 Covariance blocks}
\[
\begin{aligned}
S_{\beta\beta}&=\Lambda_{\beta|u}^{-1},\\
S_{\beta u}&=-\Lambda_{\beta|u}^{-1}B_{\beta u}D_u^{-1},\\
S_{uu}&=D_u^{-1}+D_u^{-1}B_{\beta u}^\top
\Lambda_{\beta|u}^{-1}B_{\beta u}D_u^{-1}.
\end{aligned}
\]
\textbf{English explanation.}
\begin{enumerate}[leftmargin=2em]
\item The beta covariance is the inverse Schur complement.
\item The beta-u covariance is recovered implicitly.
\item The u-u covariance includes the base inverse \(D_u^{-1}\) plus a coupling correction through beta.
\item Route B evaluates needed quadratic forms without materializing the full covariance.
\end{enumerate}

\subsection{14. Kronecker-Aware Implementation}
\paragraph{Formula 14.1 \(A^\top A\) factorization}
\[
A_n=T_n\otimes C,
\qquad
A_n^\top A_n=(T_n^\top T_n)\otimes(C^\top C),
\qquad
G:=C^\top C.
\]
\textbf{English explanation.}
\begin{enumerate}[leftmargin=2em]
\item \(A_n\) factorizes into temporal and spatial projection matrices.
\item Kronecker algebra gives the factorization of \(A_n^\top A_n\).
\item \(G=C^\top C\) is a fixed spatial Gram factor.
\item Online updates only need to refresh temporal statistics.
\end{enumerate}

\paragraph{Formula 14.2 Temporal-only transfer operator}
\[
K_{nn}=K_{nn}^{(t)}\otimes K_s,
\qquad
K_{on}=K_{on}^{(t)}\otimes K_s,
\]
\[
L_{on}=K_{on}K_{nn}^{-1}
=L_{on}^{(t)}\otimes I_s,
\qquad
L_{on}^{(t)}=K_{on}^{(t)}(K_{nn}^{(t)})^{-1}.
\]
\textbf{English explanation.}
\begin{enumerate}[leftmargin=2em]
\item Fixed spatial inducing locations make old-new and new-new covariance share the same spatial factor.
\item The spatial factor cancels in the conditional mean operator.
\item Therefore the transfer acts only along the temporal inducing dimension.
\item This is why \(L_{on}=L_{on}^{(t)}\otimes I_s\).
\end{enumerate}

\paragraph{Formula 14.3 Kronecker-preserving \(u\)-precision transfer}
\[
R_{uu,o}=B_o\otimes G,
\]
\[
R_{uu,o\rightarrow n}
=L_{on}^\top R_{uu,o}L_{on}
=\left[(L_{on}^{(t)})^\top B_oL_{on}^{(t)}\right]\otimes G.
\]
\textbf{English explanation.}
\begin{enumerate}[leftmargin=2em]
\item The old u-u likelihood precision is stored as a temporal factor times a fixed spatial factor.
\item Since \(L_{on}\) acts only temporally, the spatial factor remains \(G\).
\item The temporal factor transforms as \(L_t^\top B_oL_t\).
\item The transfer preserves the single-Kronecker structure.
\end{enumerate}

\paragraph{Formula 14.4 Updated temporal likelihood statistic}
\[
B_n=
B_{o\rightarrow n}
+
\frac1{\sigma^2}T_n^\top T_n.
\]
\[
\Lambda_{uu,n}
=
(K_{nn}^{(t)})^{-1}\otimes K_s^{-1}
+
B_n\otimes G.
\]
\textbf{English explanation.}
\begin{enumerate}[leftmargin=2em]
\item \(B_{o\rightarrow n}\) is the transferred historical temporal statistic.
\item The current block contributes \(T_n^\top T_n/\sigma^2\).
\item Their sum gives the new temporal likelihood statistic.
\item The posterior u precision is prior precision plus likelihood precision.
\item This preserves the Sylvester-compatible structure.
\end{enumerate}

\subsection{15. Structured Beta-U Cross Block}
\paragraph{Formula 15.1 Cross-block update}
\[
R_{\beta u,o\rightarrow n}
=R_{\beta u,o}(L_{on}^{(t)}\otimes I_s),
\]
\[
R_{\beta u,n}
=
R_{\beta u,o}(L_{on}^{(t)}\otimes I_s)
+
\frac1{\sigma^2}\Phi_n^\top A_n.
\]
\textbf{English explanation.}
\begin{enumerate}[leftmargin=2em]
\item \(R_{\beta u,o}\) stores historical coupling between beta and inducing variables.
\item It must be mapped into the new inducing coordinates.
\item The current block contributes new coupling information.
\item The updated cross block is the sum of transferred old coupling and new coupling.
\item This is a precision cross block, not the posterior covariance itself.
\end{enumerate}

\paragraph{Formula 15.2 Posterior cross covariance}
\[
S_{\beta u}
=
-\Lambda_{\beta|u}^{-1}B_{\beta u}D_u^{-1}.
\]
\textbf{English explanation.}
\begin{enumerate}[leftmargin=2em]
\item \(B_{\beta u}\) is a cross block in precision space.
\item The posterior cross covariance appears after inverting the full precision matrix.
\item The Schur identity gives its implicit form.
\item Route B uses it through solves and small matrix operations.
\end{enumerate}

\subsection{16. Sylvester Solves}
\paragraph{Formula 16.1 Vector-to-matrix reshape}
\[
q=\operatorname{vec}(Q),
\qquad
z=D_u^{-1}q=\operatorname{vec}(Z),
\qquad
Q,Z\in\mathbb R^{M_s\times M_t}.
\]
\textbf{English explanation.}
\begin{enumerate}[leftmargin=2em]
\item The right-hand side \(q\) is a vector of length \(M_tM_s\).
\item It is reshaped into an \(M_s\times M_t\) matrix.
\item The solution vector is reshaped in the same way.
\item This converts a Kronecker linear system into a matrix equation.
\end{enumerate}

\paragraph{Formula 16.2 Sylvester equation}
\[
D_u=(K_{nn}^{(t)})^{-1}\otimes K_s^{-1}+B_n\otimes G,
\]
\[
K_s^{-1}Z(K_{nn}^{(t)})^{-1}+GZB_n=Q.
\]
\textbf{English explanation.}
\begin{enumerate}[leftmargin=2em]
\item \(D_u\) is a sum of a prior Kronecker term and a likelihood Kronecker term.
\item The vectorization identity converts each Kronecker product into left-right matrix multiplication.
\item The prior term becomes \(K_s^{-1}ZK_t^{-1}\).
\item The likelihood term becomes \(GZB_n\).
\item The original large linear system becomes a Sylvester-type equation.
\end{enumerate}

\subsection{17. Prediction and Uncertainty}
\paragraph{Formula 17.1 Test-point features}
\[
\phi_*:=\phi(t_*,s_*),
\qquad
a_*^\top:=K_{*u}K_{uu}^{-1}.
\]
\textbf{English explanation.}
\begin{enumerate}[leftmargin=2em]
\item \(\phi_*\) is the linear feature vector for the test input.
\item \(a_*^\top\) is the sparse GP projection from inducing variables to the test point.
\item It is the single-point analogue of \(A_n\).
\end{enumerate}

\paragraph{Formula 17.2 Predictive mean}
\[
\mathbb E[y_*]
=
\phi_*^\top m_\beta+a_*^\top m_u.
\]
\textbf{English explanation.}
\begin{enumerate}[leftmargin=2em]
\item The predictive mean is the sum of the linear mean and GP residual mean.
\item \(\phi_*^\top m_\beta\) is the linear trend prediction.
\item \(a_*^\top m_u\) is the GP residual prediction.
\item The prediction remains joint.
\end{enumerate}

\paragraph{Formula 17.3 Sparse conditional residual variance}
\[
\nu_*:=k(x_*,x_*)-K_{*u}K_{uu}^{-1}K_{u*}.
\]
\textbf{English explanation.}
\begin{enumerate}[leftmargin=2em]
\item \(k(x_*,x_*)\) is the prior variance at the test point.
\item \(K_{*u}K_{uu}^{-1}K_{u*}\) is the variance explained by inducing variables.
\item The difference is the residual conditional variance of the sparse GP approximation.
\item Non-unit kernel amplitude must be included here.
\end{enumerate}

\paragraph{Formula 17.4 Structured predictive quadratic}
\[
D_uv_*=a_*,
\]
\[
\begin{bmatrix}\phi_*\\a_*\end{bmatrix}^\top
\Lambda^{-1}
\begin{bmatrix}\phi_*\\a_*\end{bmatrix}
=
a_*^\top v_*
+
\left(\phi_*-B_{\beta u}v_*\right)^\top
\Lambda_{\beta|u}^{-1}
\left(\phi_*-B_{\beta u}v_*\right).
\]
\textbf{English explanation.}
\begin{enumerate}[leftmargin=2em]
\item First solve \(v_*=D_u^{-1}a_*\) without forming \(D_u^{-1}\).
\item The left side is the posterior parameter uncertainty contribution.
\item \(a_*^\top v_*\) is the base inducing uncertainty contribution.
\item \(\phi_*-B_{\beta u}v_*\) is the effective beta feature after accounting for coupling.
\item \(\Lambda_{\beta|u}^{-1}\) is the Schur beta covariance.
\item The formula evaluates a full covariance quadratic using structured solves.
\end{enumerate}

\paragraph{Formula 17.5 Predictive variance}
\[
\operatorname{Var}(y_*)
=
\sigma^2+\nu_*+a_*^\top v_*
+
\left(\phi_*-B_{\beta u}v_*\right)^\top
\Lambda_{\beta|u}^{-1}
\left(\phi_*-B_{\beta u}v_*\right).
\]
\textbf{English explanation.}
\begin{enumerate}[leftmargin=2em]
\item \(\sigma^2\) is observation noise and is included for noisy predictive observations.
\item \(\nu_*\) is sparse conditional residual variance.
\item \(a_*^\top v_*\) is inducing posterior uncertainty.
\item The last term is the Schur beta uncertainty with beta-u coupling.
\item This is the Route B predictive variance.
\end{enumerate}

\paragraph{Formula 17.6 Mean-field comparison}
\[
\operatorname{Var}_{MF}(y_*)
=
\sigma^2+\nu_*+
\phi_*^\top S_{\beta,MF}\phi_*
+
a_*^\top S_{u,MF}a_*.
\]
\textbf{English explanation.}
\begin{enumerate}[leftmargin=2em]
\item Mean-field assumes posterior independence between \(\beta\) and \(u\).
\item The variance is just the sum of linear uncertainty and GP uncertainty.
\item It drops the beta-u cross-covariance term.
\item It is a useful ablation but not the full structured Route B predictive distribution.
\end{enumerate}

\subsection{18. Main Logical Chain}
\paragraph{Formula 18.1 Summary implication chain}
\[
\begin{aligned}
&\text{two-stage learning has stale residuals}\\
\Rightarrow\;&\text{joint training fixes the observation target}\\
\Rightarrow\;&\text{changing HiPPO bases require old-information transfer}\\
\Rightarrow\;&\text{SSGP-style old-likelihood-ratio transfer gives a principled streaming update}\\
\Rightarrow\;&\text{Route B keeps structured }\beta\text{--}u\text{ covariance}\\
\Rightarrow\;&\text{Kronecker/Sylvester structure keeps the method scalable.}
\end{aligned}
\]
\textbf{English explanation.}
\begin{enumerate}[leftmargin=2em]
\item Two-stage residual learning suffers from stale residuals.
\item Joint training models raw observations and fixes the target.
\item Changing HiPPO bases require transferring old information across coordinates.
\item Old-likelihood-ratio transfer gives a principled streaming update.
\item Route B retains beta-u structured covariance.
\item Kronecker/Sylvester algebra keeps the method scalable.
\end{enumerate}

\subsection{19. Appendix Proof: \(L_{on}=L_{on}^{(t)}\otimes I_s\)}
\paragraph{Formula 19.1 Old-new covariance factorization}
\[
\begin{aligned}
\operatorname{cov}(u_{\ell,j}^{old},u_{m,j'}^{new})
&=
\int_0^{T_o}\int_0^{T_n}
g_\ell^{(T_o)}(x)g_m^{(T_n)}(x')
k((x,z_j),(x',z_{j'}))\,dx'\,dx\\
&=
\left[
\int_0^{T_o}\int_0^{T_n}
g_\ell^{(T_o)}(x)g_m^{(T_n)}(x')k_t(x,x')\,dx'\,dx
\right]
k_s(z_j,z_{j'}).
\end{aligned}
\]
\textbf{English explanation.}
\begin{enumerate}[leftmargin=2em]
\item The covariance between old and new inducing variables is a double integral of the GP kernel.
\item The kernel inside the integral is the covariance of the latent function values.
\item Separability splits it into temporal and spatial factors.
\item The spatial factor is constant with respect to the time integrals.
\item Thus old-new covariance factorizes into temporal covariance times spatial covariance.
\end{enumerate}

\paragraph{Formula 19.2 Conditional operator proof}
\[
K_{on}=K_{on}^{(t)}\otimes K_s,
\qquad
K_{nn}=K_{nn}^{(t)}\otimes K_s.
\]
\[
\begin{aligned}
L_{on}
&=K_{on}K_{nn}^{-1}\\
&=(K_{on}^{(t)}\otimes K_s)
\left[(K_{nn}^{(t)})^{-1}\otimes K_s^{-1}\right]\\
&=K_{on}^{(t)}(K_{nn}^{(t)})^{-1}\otimes K_sK_s^{-1}\\
&=L_{on}^{(t)}\otimes I_s.
\end{aligned}
\]
\textbf{English explanation.}
\begin{enumerate}[leftmargin=2em]
\item Both covariance matrices share the same spatial factor.
\item The conditional operator is \(K_{on}K_{nn}^{-1}\).
\item Use the Kronecker inverse identity.
\item Use the Kronecker multiplication identity.
\item The spatial covariance cancels to identity.
\item Therefore the transfer is temporal-only.
\end{enumerate}

\subsection{20. Appendix Proof: Why \(R_{uu,o}=B_o\otimes G\)}
\paragraph{Formula 20.1 Likelihood precision from one block}
\[
\Lambda_{\mathrm{like},j}
=
\frac1{\sigma^2}A_j^\top A_j,
\qquad
A_j=T_j\otimes C.
\]
\[
\Lambda_{\mathrm{like},j}
=
\frac1{\sigma^2}(T_j^\top T_j)\otimes(C^\top C).
\]
\textbf{English explanation.}
\begin{enumerate}[leftmargin=2em]
\item A Gaussian likelihood contributes \(A_j^\top A_j/\sigma^2\) to inducing precision.
\item Substitute \(A_j=T_j\otimes C\).
\item Kronecker algebra separates temporal and spatial factors.
\item Every block shares the same spatial factor \(G=C^\top C\).
\end{enumerate}

\paragraph{Formula 20.2 Historical sum}
\[
R_{uu,o}
:=
\sum_{j<n}\Lambda_{\mathrm{like},j}
=
\left[
\sum_{j<n}\frac1{\sigma^2}T_j^\top T_j
\right]\otimes G.
\]
\[
B_o:=
\sum_{j<n}\frac1{\sigma^2}T_j^\top T_j,
\qquad
R_{uu,o}=B_o\otimes G.
\]
\textbf{English explanation.}
\begin{enumerate}[leftmargin=2em]
\item \(R_{uu,o}\) is the sum of historical likelihood precision contributions.
\item Since every term shares \(G\), the spatial factor can be factored out.
\item The temporal sum is called \(B_o\).
\item Hence \(R_{uu,o}=B_o\otimes G\).
\item This is an algorithmic natural-parameter invariant, not a property of arbitrary dense covariance.
\end{enumerate}

\subsection{21. Appendix: Explicit Sylvester Solver}
\paragraph{Formula 21.1 General and symmetric Sylvester forms}
\[
K_s^{-1}ZK_t^{-T}+GZB_n^\top=Q.
\]
\[
K_s^{-1}ZK_t^{-1}+GZB_n=Q
\quad\text{when matrices are symmetric.}
\]
\textbf{English explanation.}
\begin{enumerate}[leftmargin=2em]
\item The first equation is the general matrix equation.
\item Symmetry removes the transposes.
\item The implementation uses symmetrization and jitter to make this stable.
\end{enumerate}

\paragraph{Formula 21.2 Whitening}
\[
\widetilde Z=K_s^{-1/2}ZK_t^{-1/2},
\qquad
\widetilde Q=K_s^{1/2}QK_t^{1/2},
\]
\[
\widetilde G=K_s^{1/2}GK_s^{1/2},
\qquad
\widetilde B=K_t^{1/2}B_nK_t^{1/2}.
\]
\[
\widetilde Z+\widetilde G\widetilde Z\widetilde B=\widetilde Q.
\]
\textbf{English explanation.}
\begin{enumerate}[leftmargin=2em]
\item Whitening turns the prior precision part into an identity term.
\item \(\widetilde Z\) is the whitened solution matrix.
\item \(\widetilde Q\) is the transformed right-hand side.
\item \(\widetilde G\) and \(\widetilde B\) are whitened likelihood factors.
\item The resulting equation is easier to diagonalize.
\end{enumerate}

\paragraph{Formula 21.3 Diagonal solution}
\[
\widetilde G=U_s\Gamma U_s^\top,
\qquad
\widetilde B=U_tHU_t^\top.
\]
\[
\widehat Z=U_s^\top\widetilde ZU_t,
\qquad
\widehat Q=U_s^\top\widetilde QU_t.
\]
\[
(1+\gamma_i\eta_j)\widehat Z_{ij}=\widehat Q_{ij},
\qquad
\widehat Z_{ij}=\frac{\widehat Q_{ij}}{1+\gamma_i\eta_j}.
\]
\textbf{English explanation.}
\begin{enumerate}[leftmargin=2em]
\item Eigendecompose the whitened spatial and temporal likelihood factors.
\item Rotate the solution and right-hand side into the eigenbasis.
\item Both factors become diagonal.
\item The matrix equation decouples entry by entry.
\item Each entry is solved by a scalar division.
\end{enumerate}

\paragraph{Formula 21.4 Fast quadratic form}
\[
q^\top D_u^{-1}q
=
\sum_{i=1}^{M_s}\sum_{j=1}^{M_t}
\frac{\widehat Q_{ij}^2}{1+\gamma_i\eta_j}.
\]
\textbf{English explanation.}
\begin{enumerate}[leftmargin=2em]
\item Prediction often needs only the scalar quadratic form.
\item In diagonal coordinates, the solution entry is \(\widehat Q_{ij}/(1+\gamma_i\eta_j)\).
\item The inner product is the sum of \(\widehat Q_{ij}\widehat Z_{ij}\).
\item Substitution gives the fast quadratic formula.
\item This avoids reconstructing the full solution when only uncertainty is needed.
\end{enumerate}

\paragraph{Formula 21.5 Complexity}
\[
\mathcal O(M_t^3+M_s^3+M_t^2M_s+M_tM_s^2).
\]
\textbf{English explanation.}
\begin{enumerate}[leftmargin=2em]
\item \(M_t^3\) comes from temporal eigendecomposition.
\item \(M_s^3\) comes from spatial eigendecomposition.
\item The mixed terms come from two-sided matrix transforms.
\item This is much cheaper than a dense solve over \(M_tM_s\) variables.
\item Fixed spatial geometry allows caching spatial factors.
\end{enumerate}

\subsection{22. When the Single-Kronecker Assumption Can Fail}
\paragraph{Formula 22.1 Changing spatial pattern}
\[
C_j\neq C,
\qquad
G_j=C_j^\top C_j,
\qquad
R_{uu,o}=\sum_{j<n}B_j\otimes G_j.
\]
\textbf{English explanation.}
\begin{enumerate}[leftmargin=2em]
\item If spatial observation patterns change, the spatial projection can become block-dependent.
\item Then the spatial Gram factor also changes.
\item The historical precision is no longer a single Kronecker product.
\item It becomes a sum of Kronecker products.
\end{enumerate}

\paragraph{Formula 22.2 Sum-of-Kronecker extension}
\[
R_{uu,o}=\sum_{r=1}^{R_o}B_r\otimes G_r.
\]
\[
D_u=K_t^{-1}\otimes K_s^{-1}+\sum_{r=1}^R B_r\otimes G_r.
\]
\[
K_s^{-1}MK_t^{-1}+\sum_{r=1}^RG_rMB_r=H.
\]
\textbf{English explanation.}
\begin{enumerate}[leftmargin=2em]
\item The first equation is a more general historical likelihood representation.
\item The second equation adds it to the prior precision.
\item The third equation is the corresponding matrix equation.
\item It is no longer the simple single-term Sylvester equation.
\item This is a future extension rather than the current main Route B implementation.
\end{enumerate}


\end{document}
